\documentclass[11pt,a4paper]{article}

% ============================================================
% Standard English mathematical typesetting preamble
% ============================================================
\usepackage[utf8]{inputenc}
\usepackage[T1]{fontenc}
\usepackage[english]{babel}
\usepackage{amsmath,amssymb,amsthm}
\usepackage{geometry}
\geometry{margin=1in}
\usepackage{graphicx}
\usepackage{booktabs}
\usepackage{array}
\usepackage{fancyhdr}
\usepackage{titlesec}
\usepackage{enumitem}
\usepackage{mathtools}
\usepackage{bm}

% Theorem-like environments
\theoremstyle{plain}
\newtheorem{theorem}{Theorem}
\newtheorem{lemma}{Lemma}
\theoremstyle{definition}
\newtheorem{remark}{Remark}

% Page style
\pagestyle{fancy}
\fancyhf{}
\fancyhead[C]{\small Cayley --- Collected Mathematical Papers, Vol.\ VIII}
\fancyfoot[C]{\thepage}
\renewcommand{\headrulewidth}{0.4pt}

% Section formatting to match original numbered articles
\titleformat{\section}{\large\scshape\centering}{}{0em}{}
\titleformat{\subsection}{\normalsize\bfseries}{}{0em}{}

% Custom commands for Cayley's notation
\newcommand{\dd}{\,d}
\DeclareMathOperator{\Jac}{Jac}

\begin{document}

% ============================================================
% PAPER 520 (continued from earlier pages; here covers book pp.~317--365)
% ============================================================
\section{[520]}
\subsection{On the Centro-Surface of an Ellipsoid (continued).}

\begin{center}
\textit{[Continuation of the memoir; book pages 317--365.]}
\end{center}

% --- p.317 ---
\bigskip
\textit{p.~317 [520]}
\smallskip

\noindent
imaginary is of course attained by the equation. The new results suitably modified
would be applicable to the theory treated from the second point of view; but I do not
on the present occasion attempt so to present them.

\subsubsection*{The Ellipsoid; Parameters $\xi,\eta,\zeta$. Art.\ Nos.\ 1--6.}

\textbf{1.} The position of a point $(X, Y, Z)$ on the ellipsoid
\[
\frac{X^2}{a} + \frac{Y^2}{b} + \frac{Z^2}{c} = 1
\]
may be determined by means of the parameters, or elliptic coordinates, $\xi, \eta$: viz.\ these are such that we have
\[
\frac{X^2}{a + \xi} + \frac{Y^2}{b + \xi} + \frac{Z^2}{c + \xi} = 1,\qquad
\frac{X^2}{a + \eta} + \frac{Y^2}{b + \eta} + \frac{Z^2}{c + \eta} = 1,
\]
or, what is the same thing, $\xi, \eta$ are the roots of the quadric equation
\[
\frac{X^2}{a + \theta} + \frac{Y^2}{b + \theta} + \frac{Z^2}{c + \theta} = 1.
\]

(In its actual form this is a cubic equation, but there is a root $\theta = 0$, which is to be thrown out, and the quadric equation is
\[
\theta^2 + a(b + c + b + c - X^2 - Y^2 - Z^2)
\]
\[
+ \{bc + ca + ab - (b + c)X^2 - (c + a)Y^2 - (a + b)Z^2\} = 0,
\]
or putting
\[
P = a + b + c,\qquad Q = bc + ca + ab,\qquad R = abc,
\]
the equation is
\[
\theta^2 + (P - X^2 - Y^2 - Z^2)\theta + Q - (b + c)X^2 - (c + a)Y^2 - (a + b)Z^2 = 0.)
\]

\textbf{2.} It is convenient to write throughout
\begin{align*}
P' &= b + c,\\
Q' &= bc,\\
P'' &= c + a,\\
Q'' &= ca,\\
P''' &= a + b,\\
Q''' &= ab.
\end{align*}

(whence $a + \beta + \gamma = 0$),

\clearpage

% --- p.318 ---
\bigskip
\textit{p.~318 [520]}
\smallskip

As usual, $a$ is taken to be the greatest and $c$ the least of the semi-axes; we have thus $a, \eta$ each of them positive, and $\beta$ negative, $=-b'$ where $b'$ is a positive quantity $=b - c$. A distinction arises in the sequel between the two cases $a + \beta = 2b'$ and $a + \beta < 2b'$; but the two cases are not essentially different, and it is convenient to assume $a + \beta > 2b'$, that is, $a > b - 2b' + c$ or $\gamma > a$, say $\gamma > 0$ positive. The limiting case $a + \beta = 2b'$ or $\gamma = a$ requires special consideration.

\textbf{3.} We have
\begin{align*}
-\beta\gamma X^2 &= a(a^2 + \xi)(a^2 + \eta),\\
-\gamma a Y^2 &= b(b^2 + \xi)(b^2 + \eta),\\
-a\beta Z^2 &= c(c^2 + \xi)(c^2 + \eta).
\end{align*}

It is in fact easy to verify that these values satisfy as well the equation of the ellipsoid as the assumed equations defining the elliptic coordinates $\xi, \eta$. We may also obtain the relations
\[
X^2 + Y^2 + Z^2 = a + b + c + \xi + \eta,
\]
\[
aX^2 + bY^2 + cZ^2 = a^2 + b^2 + c^2 + (b + c)\xi + (c + a)\eta + (a + b)\zeta,
\]
\[
\tfrac{1}{a}X^2 + \tfrac{1}{b}Y^2 + \tfrac{1}{c}Z^2 = 1 + \xi + \eta + \tfrac{\xi\eta}{R}.
\]

These, however, are obtained more readily from the equation in $a$, viz.\ the roots thereof being $\xi, \eta$, we have
\[
-\xi - \eta = P + b + c - X^2 - Y^2 - Z^2,
\]
\[
\xi\eta = bc + ca + ab - (b + c)X^2 - (c + a)Y^2 - (a + b)Z^2,
\]
which lead at once to the relations in question.

\textbf{4.} Considering $\xi$ as constant, the locus of the point $(X, Y, Z)$ is the intersection of the ellipsoid with the confocal ellipsoid
\[
\frac{X^2}{a + \xi} + \frac{Y^2}{b + \xi} + \frac{Z^2}{c + \xi} = 1;
\]
viz.\ this is one of the curves of curvature through the point; and similarly considering $\eta$ as constant, the locus of the point is the intersection of the ellipsoid with the confocal ellipsoid
\[
\frac{X^2}{a + \eta} + \frac{Y^2}{b + \eta} + \frac{Z^2}{c + \eta} = 1;
\]
viz.\ this is the other of the curves of curvature through the point.

\textbf{5.} If instead of $\xi$ and $\eta$ we write $k$ and $l$, we may consider $k$ as extending between the values $-a^2, -b^2$, and $l$ as extending between the values $-b^2, -c^2$.

A root will thus give the series of curves of curvature one of which is the section by the plane $X = 0$, or ellipse semi-axes $b, c$; say this is the minor-mean-axes section. In particular $k = -a^2$ gives the ellipse just referred to; and $k = -b^2$, or say $k = -b^2 + \epsilon$, gives two indefinitely near curves of curvature, viz.\ a curve of curvature portions extends from an umbilicus above the plane of $xy$, through the extremity of the semi-axis $a$ to an umbilicus above the plane of $xy$.

The ellipse last referred to may be called the multilinear section; the other two principal sections being the major-mean section and the minor-mean section respectively.

In the limiting case $k = l = -b^2$, we have the umbilici, viz.\ these are given by
\[
\frac{X^2}{a - b} = \frac{Z^2}{b - c},\qquad Y = 0,\qquad \frac{X^2}{a} + \frac{Z^2}{c} = 1.
\]

The two series of curves of curvature cover the whole real surface of the ellipsoid; so that at any real point thereof we have $\xi = k, \eta = l$; or else if $\xi, \eta$ are negative real values lying within the foregoing limits $-a^2, -b^2$ and $-b^2, -c^2$ respectively. But observe that $\xi, \eta$ taken separately may each extend between the limits $-a^2, -c^2$.

\textbf{6.} Suppose $\xi = \eta$, the equation is a will have equal roots, or the condition is
\[
(P - X^2 - Y^2 - Z^2)^2 = 4\{Q - (b + c)X^2 - (c + a)Y^2 - (a + b)Z^2\};
\]
viz.\ the surface by its intersection with the ellipsoid determines the envelope of the curves of curvature. This envelope is in fact a system of eight imaginary lines, four of them belonging to one of the systems of right lines on the ellipsoid, the other four to the other of the systems of right lines on the ellipsoid. For in the values of $X, Y, Z$ writing $\eta = \xi$, we find
\begin{align*}
\tfrac{X}{a}\sqrt{-\beta\gamma} &= a + \xi,\\
\tfrac{Y}{b}\sqrt{-\gamma a} &= b + \xi,\\
\tfrac{Z}{c}\sqrt{-a\beta} &= c + \xi,
\end{align*}
or representing for shortness the left-hand functions by $\tfrac{X}{a}\sqrt{\phantom{x}}, \tfrac{Y}{b}\sqrt{\phantom{x}}, \tfrac{Z}{c}\sqrt{\phantom{x}}$, the right lines are
\begin{align*}
a + \xi &= X,& b + \xi &= Y,& c + \xi &= Z,\\
a + \xi &= X,& b + \xi &= -Y,& c + \xi &= -Z,\\
a + \xi &= -X,& b + \xi &= Y,& c + \xi &= -Z,\\
a + \xi &= -X,& b + \xi &= -Y,& c + \xi &= Z.
\end{align*}

\clearpage

% --- p.320 ---
\bigskip
\textit{p.~320 [520]}
\smallskip

so that in the two trends each line intersects the four lines of the other trend, but it does not intersect the remaining three lines of its own trend. The intersections are four points corresponding to $\xi = -a^2$, being the imaginary umbilici in the plane $X = 0$: four to $\xi = -b^2$, being the real umbilici in the plane $Y = 0$: four to $\xi = -c^2$, being the imaginary umbilici in the plane $Z = 0$: and four corresponding to $\xi = a\beta\gamma$, which may be called the umbilici at infinity ($1$).

\subsubsection*{Separable and Concomitant Centro-curves. Art.\ No.\ 7.}

\textbf{7.} Consider any particular curve of curvature; the normals at the several points thereof successively intersect each other in a series of points forming a curve; and we have thus, corresponding to the particular curve of curvature, a curve on the centro-surface, which curve may be called the separated centro-curve. Again the same normals, viz.\ those at the several points of the particular curve of curvature, are intersected by the normal at each point by the consecutive normal belonging to the other curve of curvature through that point; and we have thus, corresponding to the particular curve of curvature, a curve on the centro-surface, which curve may be called the concomitant centro-curve. If instead of a single curve of curvature we consider the whole series, say the major-mean curves of curvature, we have a series of major-mean separated centro-curves, and also a series of major-mean concomitant centro-curves; and similarly considering the series of the minor-mean curves of curvature, we have a series of minor-mean separated centro-curves, and also a series of minor-mean concomitant centro-curves; the conjunction of the several curves will be discussed further on; but it may be benevolent to remark here that the centro-surface may be considered as consisting of two portions, say,

(A) locus of the major-mean separated centro-curves; and also of the minor-mean concomitant centro-curves;

(B) locus of the minor-mean separated centro-curves, and also of the major-mean concomitant centro-curves.

\subsubsection*{Investigation of expressions for the Coordinates of a point on the Centro-surface. Art.\ Nos.\ 8 to 13.}

\textbf{8.} Consider the normal at the point $(X, Y, Z)$. Taking in the first instance $(x, y, z)$ as current coordinates, the equations are
\[
\frac{x - X}{\tfrac{X}{a}} = \frac{y - Y}{\tfrac{Y}{b}} = \frac{z - Z}{\tfrac{Z}{c}} = \lambda \ \text{suppose}.
\]

\smallskip
\footnotesize
$^1$ According to Salmon, \textit{Solid Geometry,} [Ed.\ 2, 1865], p.~210, the number of umbilici is $a$ surface of the $n$th order is $= n(10n^2 - 15n + 6)$, viz.\ for $n = 2$, $= 14$, or it is, in the ordinary theory, and recognising the umbilici at infinity. For surfaces properly umbilici or not, the 4 points which I call the umbilici at infinity in the present theory dimension themselves in like manner with the 12 ordinary umbilici.

\clearpage

% --- p.321 ---
\bigskip
\textit{p.~321 [520]}
\smallskip

or, what is the same thing,
\[
x = X\left(1 + \tfrac{\lambda}{a}\right),\quad y = Y\left(1 + \tfrac{\lambda}{b}\right),\quad z = Z\left(1 + \tfrac{\lambda}{c}\right).
\]

Suppose now that the normal meets the consecutive normal, or normal at the point $X + dX, Y + dY, Z + dZ$; and let $x, y, z$ belong to the point of intersection of the two normals; we have
\begin{align*}
0 &= dX\left(1 + \tfrac{\lambda}{a}\right) + \tfrac{X}{a}\,d\lambda,\\
0 &= dY\left(1 + \tfrac{\lambda}{b}\right) + \tfrac{Y}{b}\,d\lambda,\\
0 &= dZ\left(1 + \tfrac{\lambda}{c}\right) + \tfrac{Z}{c}\,d\lambda;
\end{align*}
which determine the direction of the consecutive point; the equations in fact give
\[
0 = \begin{vmatrix} dX, & \tfrac{X}{a},\\ dY, & \tfrac{Y}{b},\\ dZ, & \tfrac{Z}{c} \end{vmatrix},
\]
or, what is the same thing,
\[
0 = \begin{vmatrix} dX, & dY, & dZ,\\ b cX, & caY, & abZ \end{vmatrix},
\]
which is the differential equation of the curve of curvature. This equation must therefore be satisfied by taking for $X + dX, Y + dZ, Z + dZ$ the coordinates of the consecutive point along either of the curves of curvature,---say along that which is the intersection with the surface
\[
\frac{X^2}{a^2 + \eta} + \frac{Y^2}{b^2 + \eta} + \frac{Z^2}{c^2 + \eta} = 1.
\]

\textbf{9.} To verify this, observe that we then have
\[
\frac{X dX}{a^2} + \frac{Y dY}{b^2} + \frac{Z dZ}{c^2} = 0,
\]
\[
\frac{X dX}{a^2 + \eta} + \frac{Y dY}{b^2 + \eta} + \frac{Z dZ}{c^2 + \eta} = 0,
\]
or, what is the same thing,
\[
X dX : Y dY : Z dZ = a^2(a^2 + \eta)\beta : b^2(b^2 + \eta)\gamma : c^2(c^2 + \eta)a,
\]

\clearpage

% --- p.322 ---
\bigskip
\textit{p.~322 [520]}
\smallskip

But from the equations $-\beta\gamma X^2 = a^2(a^2 + \xi)(a^2 + \eta)$, etc., these become
\[
X dX : Y dY : Z dZ = \tfrac{X^2}{a^2 + \xi} : \tfrac{Y^2}{b^2 + \xi} : \tfrac{Z^2}{c^2 + \xi};
\]
or, what is the same thing,
\[
dX : dY : dZ = \tfrac{X}{a^2 + \xi} : \tfrac{Y}{b^2 + \xi} : \tfrac{Z}{c^2 + \xi};
\]
and, substituting these values in the determinant equation, it becomes
\[
0 = \begin{vmatrix} \tfrac{X}{a^2 + \xi} & \tfrac{Y}{b^2 + \xi} & \tfrac{Z}{c^2 + \xi}\\
\tfrac{X}{a} & \tfrac{Y}{b} & \tfrac{Z}{c}\\
b cX, & caY, & abZ \end{vmatrix},
\]
which is identically true, since evidently the determinant vanishes.

\textbf{10.} Proceeding with the solution, we have from the three equations
\[
X dX + Y dY + Z dZ = d\lambda\!\left(\tfrac{X^2}{a^2} + \tfrac{Y^2}{b^2} + \tfrac{Z^2}{c^2}\right) + d\lambda\!\left(\tfrac{X^2}{a^2}\,\tfrac{1}{a^2} + \tfrac{Y^2}{b^2}\,\tfrac{1}{b^2} + \tfrac{Z^2}{c^2}\,\tfrac{1}{c^2}\right) = 0,
\]
and observing that from the equation
\[
X^2 + Y^2 + Z^2 = a + b + c + \xi + \eta,
\]
considering therein $\eta$ as constant, we have
\[
X dX + Y dY + Z dZ = \tfrac{1}{2}d\xi,
\]
the equation becomes
\[
\tfrac{1}{2}d\xi + d\lambda = 0;
\]
and the three equations then are
\[
0 = dX\left(1 + \tfrac{\lambda}{a}\right) - \tfrac{X}{2 a}\,d\xi, \text{ etc.,}
\]
or say
\[
0 = dX(a + \lambda) - \tfrac{1}{2}X d\xi, \text{ etc.}
\]

But from the equation $-\beta\gamma X^2 = a^2(a^2 + \xi)(a^2 + \eta)$, considering therein $\eta$ as a constant, we have
\[
\frac{dX}{X} = \frac{1}{2}\frac{d\xi}{a + \xi},
\]
and the equations then become
\[
0 = \tfrac{a + \lambda}{a + \xi} - 1, \text{ etc.,}
\]
viz.\ these are all satisfied if only $\lambda = \xi$.

\textbf{11.} The coordinates of the point of intersection of the two normals thus are
\[
x = X\left(1 + \tfrac{\xi}{a}\right),\quad y = Y\left(1 + \tfrac{\xi}{b}\right),\quad z = Z\left(1 + \tfrac{\xi}{c}\right).
\]

\clearpage

% --- p.323 ---
\bigskip
\textit{p.~323 [520]}
\smallskip

or squaring, and substituting for $X^2$, etc., their values as given by
\[
-\beta\gamma X^2 = a^2(a^2 + \xi)(a^2 + \eta), \text{ etc.,}
\]
the equations become
\begin{align*}
-\beta\gamma x^2 &= a^2(a + \xi)^3(a + \eta), \text{ etc.,}\\
-\gamma a y^2 &= b(b + \xi)^3(b + \eta), \text{ etc.,}\\
-a\beta z^2 &= c(c + \xi)^3(c + \eta), \text{ etc.;}
\end{align*}
viz.\ these equations give $(x, y, z)$ the coordinates of a point on the centro-surface, the intersection of the normal at the point $(X, Y, Z)$ of the ellipsoid, (determined by the parameters $\xi, \eta$), by the normal at the consecutive point along the curve of curvature
\[
\frac{X^2}{a^2 + \eta} + \frac{Y^2}{b^2 + \eta} + \frac{Z^2}{c^2 + \eta} = 1.
\]

Of course by interchanging $\xi$ and $\eta$, we should obtain the coordinates of the point of intersection of the normal at the same point $(X, Y, Z)$ by the normal at the consecutive point along the other curve of curvature
\[
\frac{X^2}{a^2 + \xi} + \frac{Y^2}{b^2 + \xi} + \frac{Z^2}{c^2 + \xi} = 1.
\]

\textbf{12.} I stop for a moment to consider the foregoing two equations
\[
\lambda = \xi,\quad d\lambda = -\tfrac{1}{2}d\xi.
\]
which at first sight appear inconsistent. But observe that in the foregoing solution $\lambda$ is the parameter of the point $(x, y, z)$ of the centro-surface considered as a point on the normal at $(X, Y, Z)$; $\lambda + d\lambda$ is the parameter of the same point considered as a point on the normal at the consecutive point $(X + dX, Y + dY, Z + dZ)$; the value $\lambda + d\lambda = \xi - \tfrac{1}{2}d\xi$ would belong to a different point, viz.\ the consecutive point of the centro-surface considered as a point on the consecutive normal: whence the $d\lambda$ of the solution ought not to be $= d\xi$. In further explanation, observe that the equations
\[
x = X\!\left(1 + \tfrac{\lambda}{a}\right), \text{ etc.,} \quad\lambda = \xi,
\]
if we pass from $(x, y, z)$ to its consecutive point on the centro-surface, give
\[
dx = dX\!\left(1 + \tfrac{\lambda}{a}\right) + \tfrac{X}{a}\,d\xi,
\]
but since by the equations
\[
0 = dX\!\left(1 + \tfrac{\lambda}{a}\right) + \tfrac{X}{a}\,d\lambda,
\]
this is
\[
dx = \tfrac{X}{a}d\xi - \tfrac{X}{a}d\lambda,
\]

\smallskip
\footnotesize
$^1$ The expressions are given in effect, but not explicitly, Salmon, p.~149.

\clearpage

% --- p.324 ---
\bigskip
\textit{p.~324 [520]}
\smallskip

Or since
\[
a' = \lambda = \xi\quad(a + \xi),
\]
this is
\[
\frac{dx}{x} = \tfrac{1}{a + \xi}\,d\xi,
\]
and similarly
\[
\frac{dy}{y} = \tfrac{1}{b + \xi}\,d\xi,\quad
\frac{dz}{z} = \tfrac{1}{c + \xi}\,d\xi,
\]
which are the correct values of $dx, dy, dz$ as derived from the equations
\[
-\beta\gamma x^2 = a^2(a + \xi)^3(a + \eta), \text{ etc.,}
\]
so that taking $\eta$ at pleasure and considering $\eta$ as denoting this function of $\eta$, the values of $\xi, \eta$ belonging to a point on the nodal curve are $\xi = a^2(a + \beta), \, \xi = c(a + \beta) c$; and the value of $\eta$ is then given as before.

\textbf{13.} The form just given is analytically the most convenient, but there is some advantage in writing $\tfrac{1}{a + \xi}, \tfrac{1}{a + \eta}$ in the place of $a, \eta$ respectively; viz.\ we then have expressions for the coordinates $(x, y, z)$ of the centro-surface in terms of the coordinates $(x, y, z)$ on the centro-surface, and writing for shortness $f$ for the elimination of $(\xi, \eta)$ from these expressions will therefore lead to the equation of the centro-surface.

\textit{Discussion by means of the equations $-\beta\gamma x^2 = a^2(a + \xi)^3(a + \eta),$ etc., Principal Sections, etc. Art.\ Nos.\ 14 to 24 (general subheading).}

\textbf{14.} To fix the ideas consider the section of the surface by the plane $z = 0$; we have in the surface $z = 0$, that is, $\xi = -c^2$, or else $\eta = -c^2$, values which give respectively
\begin{align*}
-\beta\gamma x^2 &= a(a - c)^3(a - c), & -\beta\gamma x^2 &= a^2(a^2 - c^2)^3 a,\\
-\gamma a y^2 &= b(b - c)^3(b - c), & \text{or}\quad -\gamma a y^2 &= b^2(b^2 - c^2)^3 b,
\end{align*}
or, what is the same thing,
\begin{align*}
\tfrac{x^2}{(a - c)^3} &= \tfrac{a}{-\beta\gamma}(a + \eta), & \tfrac{x^2}{(a^2 - c^2)^3} &= \tfrac{a}{-\beta\gamma}(a + \xi),\\
\tfrac{y^2}{(b - c)^3} &= \tfrac{b}{-a\gamma}(b + \eta), & \tfrac{y^2}{(b^2 - c^2)^3} &= \tfrac{b}{-a\gamma}(b + \xi).
\end{align*}

The first set of equations gives
\[
\frac{a^2 x^2}{(a - c)^3} + \frac{b^2 y^2}{(b - c)^3} = 1,
\]
which is the equation of an ellipse.

The second set of equations gives
\[
(a x^2 + b y^2)^2 = 2(a^2 b^2 c^2 a y^2)(b x^2 + a y^2),
\]
which is the equation of an evolute of an ellipse.

\clearpage

% --- p.325 ---
\bigskip
\textit{p.~325 [520]}
\smallskip

\textbf{15.} The ellipse $\tfrac{a^2 x^2}{(a - c)^3} + \tfrac{b^2 y^2}{(b - c)^3} = 1$ is a cuspidal curve on the surface, and the section by the plane $z = 0$ is consequently made up of this ellipse counted three times, and of the evolute; it is therefore of the twelfth order; and the order of the surface is in fact $= 12$.

It is clear that the section of the centro-surface arises from the section $\tfrac{X^2}{a} + \tfrac{Y^2}{b} = 1$ viz.\ the normal at any point of this ellipse lies in the plane $Z = 0$; the latter section being by a normal at the consecutive point of the other curve of curvature gives a point on the ellipse $\tfrac{a^2 x^2}{(a - c)^3} + \tfrac{b^2 y^2}{(b - c)^3} = 1$, which ellipse is therefore the concomitant centro-curve. Observe that this other curve of curvature cuts the ellipse $\tfrac{X^2}{a} + \tfrac{Y^2}{b} = 1$ at right angles, and that the normals at the consecutive points about meet in the same point on the ellipse $\tfrac{a^2 x^2}{(a - c)^3} + \tfrac{b^2 y^2}{(b - c)^3} = 1$; this shows that the last-mentioned ellipse is a cuspidal curve on the centro-surface.

\textbf{16.} The three principal sections of the centro-surface are consequently
\begin{align*}
z = 0,\quad &\frac{a^2 x^2}{(a - c)^3} + \frac{b^2 y^2}{(b - c)^3} = 1,\quad\text{and}\quad (a x^2)^2 = (b x^2)^2,\\
y = 0,\quad &\frac{a^2 x^2}{(a - c)^3} + \frac{c^2 z^2}{(b - c)^3} = 1,\quad\text{and}\quad (a x^2)^2 = (c z^2)^2,\\
x = 0,\quad &\frac{b^2 y^2}{(a - c)^3} + \frac{c^2 z^2}{(b - c)^3} = 1,\quad\text{and}\quad (b y^2)^2 = (c z^2)^2;
\end{align*}
viz.\ each section is made up of an ellipse counting three times and of an evolute of an ellipse. The figure for shortness represented the three evolutes by their principal cusps it of the like character.

\textbf{17.} Considering only the positive directions of the axes, we have on each axis two points, viz.\
\begin{center}
\begin{tabular}{cccc}
axis of $x$ & $x = \tfrac{a - c}{a}$, & $x = \tfrac{a - b}{a}$;\\
axis of $y$ & $y = \tfrac{b - c}{b}$, & $y = \tfrac{a - b}{b}$;\\
axis of $z$ & $z = \tfrac{b - c}{c}$, & $z = \tfrac{a - c}{c}$;
\end{tabular}
\end{center}

\clearpage

% --- p.326 ---
\bigskip
\textit{p.~326 [520]}
\smallskip

through each of which, in the two different planes passes through the axis respectively, there passes an ellipse and an evolute. In the second case $a + \beta > 2b'$, the disposition of the points is as shown in the figure.

\begin{center}
\textit{[Figure: principal sections diagram --- omitted]}
\end{center}

Plane of $xz$, evolute is outside ellipse:
\begin{itemize}[leftmargin=2em]
\item[] $yz$, inside;
\item[] $xy$, cuts.
\end{itemize}

but in the contrary case $a + \beta < 2b'$, the disposition is

Plane of $xz$, evolute is outside ellipse:
\begin{itemize}[leftmargin=2em]
\item[] $yz$, cuts;
\item[] $xy$, inside;
\end{itemize}

there is no real difference, and to fix the ideas I attend exclusively to the first-mentioned case
\[
a + \beta > 2b'.
\]

\textbf{18.} In each of the principal planes, the evolute and ellipse, qua curves of the orders 6 and 2 respectively, intersect in twelve points, 2 in each quadrant; viz.\ of the 3 points two unite together into a twofold point or point of contact, and the third is a point of simple intersection; assuming for the moment that this is so, the figure at once shows that in the plane of $xy$ or unilinear plane the contact is real, the intersection imaginary; in the plane of $yz$, or major-mean plane, the contact is imaginary, the intersection real; but in the plane of $xz$ or minor-mean plane the contact and intersection are each imaginary. The contacts are, as will appear later, the umbilici of the centro-surface; thus the same name "umbilicus centres," or "umbilicus", the simple intersections "points of entropy," or simply "entropy." By what precedes there are real umbilici, and the case the umbilicus (or each quadrant one); and in the major-mean plane four real entropies (in each quadrant one); the other umbilicar centres and entropies are respectively imaginary.

\clearpage

% --- p.327 ---
\bigskip
\textit{p.~327 [520]}
\smallskip

\textbf{19.} The surface consists of two sheets intersecting in a nodal curve connecting the entropy with the umbilicar centres. As to the form of this curve there is a cusp at the entropy, and the curve does not terminate at the umbilicar centre but, on passing it, from crunodal becomes acnodal (viz.\ there is no longer through the curve any real sheet of the surface); moreover the curve is not at the umbilicar centre

\begin{center}
\textit{[Figure: nodal curve with cusp at entropy diagram --- omitted]}
\end{center}

perpendicular to the plane of $xy$, say there is consequently on the opposite side of the plane a symmetrically situate branch of the curve, viz.\ the umbilicar centre is a node on the nodal curve. Completing the curve, the nodal curve consists of two distinct portions; one on the positive side of the plane of $yz$ or minor-mean plane consisting of two cuspidal branches as shown in the figure; the other a symmetrically situate portion on the negative side of the minor-mean plane.

\subsubsection*{Intersections of Evolute and Ellipse.}

\textbf{20.} Consider it the plane of $xy$ the ellipse and evolute,
\[
\frac{a^2 x^2}{(a - c)^3} + \frac{b^2 y^2}{(b - c)^3} = 1,\quad (a x^2)^2 + (b y^2)^2 - c^2(a^2 y^2 + b^2 x^2) y^2 = 0.
\]

First, these are satisfied by
\[
a^2 x^2 = -\tfrac{a^2}{-\beta\gamma}(b^2 + \gamma)^3,\quad
b^2 y^2 = -\tfrac{b^2}{-\gamma a}(a^2 + \gamma)^3;
\]
viz.\ the equations respectively become
\[
\tfrac{-(\beta^2 + \gamma)^3}{-\beta\gamma} + \tfrac{(-a^2 + \gamma)^3}{-\gamma a} = 2(a + \beta)c.
\]

The first of which is $a + \beta + \alpha\gamma + \beta\gamma$, and the second is
\[
\sqrt{1}\{a^2(c - \beta)^3 + 2c\sqrt{\gamma a}\sqrt{-\beta\gamma(a - \beta + b)^3} + 4(c - \beta)^3\} = 0.
\]

But the equation in $\sqrt{1}\{a^2(c - \beta)^3 + 2c(a^3 + b^3)\sqrt{-\beta\gamma} + \cdots\} = 0$,
or, what is the same thing,
\[
(a^2 - \beta^2)(a^2 b^2 + 2a c \gamma + \cdots)^2 - 4ay^2 z^2(a + \beta)^3 + 2c(a^2 - \beta^2)\cdots = 0.
\]

\clearpage

% --- p.328 ---
\bigskip
\textit{p.~328 [520]}
\smallskip

and substituting the foregoing values, this is
\[
(a^2 - \beta^2)\{\tfrac{a^2 + \beta^2}{\gamma}\,a^2\gamma^3 - a^2 \beta^3\gamma\} + 8\gamma^2(\beta + \gamma)(a^2 + \gamma)^2 - 8a\beta\gamma^3 = 0,
\]
that is,
\[
\tfrac{1}{\gamma}\,a\beta^3(a^2 + \beta^2)(a + \beta + \gamma)(a + \beta) = 0,
\]
which, putting therein $a + \beta + \gamma = a$, and $a^2 + \beta^2 + a^2 + a^2 - 2a\beta\gamma$, is satisfied; that is, in the points on points or points of contact of the values and evolute.

\textbf{21.} Secondly, consider the values
\begin{align*}
a^2 x^2 &= \tfrac{a}{-\beta\gamma}\!\left(\tfrac{(y - a)^3}{\gamma - a}\right) (\text{Coordinates of entropies in plane of $xy$ (real).}),\\
b^2 y^2 &= \tfrac{b}{-\gamma a}\!\left(\tfrac{(\beta - a)^3}{\gamma - \beta}\right);
\end{align*}

Substituting in the equation of the ellipse, we have
\[
\tfrac{a(\beta - a)^3 + \beta\gamma a(\gamma - \beta)(a + \beta + \gamma)}{(\gamma - a)(\gamma - \beta)} = 1,
\]
which is satisfied identically; and substituting in the equation of the evolute, we have first
\[
a^2 x^2 + b^2 y^2 = -\tfrac{a^2(\beta - a)^3 + b^3(\gamma - a)(a + \beta - a)}{\gamma a(\gamma - a)} = 0,
\]
or the equation is satisfied identically: and substituting in the equation of the evolute, we have
\[
\sqrt{a^2(a - \beta)^3 \beta^3 c^2 + 2a\gamma a(\gamma - \beta)^3(a + \beta - \gamma)} - \sqrt{\gamma a}\sqrt{a^2(\gamma - \beta)^3(\gamma - a)} = 0,
\]
and thus, in virtue of $a(\beta - a) + \beta(\gamma - a) + \gamma(a + \beta) = 0$ becomes
\[
-\tfrac{3a\beta(\beta - \gamma)(\gamma - \beta) + b(\beta - a)(\gamma - a)}{\gamma a(\gamma - a)} = -\tfrac{3a(\beta - \gamma)(\gamma - a)}{a(\gamma - a)},
\]
and thus, completing the substitution, it is seen that the equation of the evolute is also satisfied. The points last considered are simple intersections, and we have thus the complete number $(3 + 4, = 12)$ of the intersections of the evolute and ellipse.

\textbf{22.} We have $a, \eta$ positive, $\beta$ negative; whence $a - \beta$ can positive, $\beta - \gamma$ negative; $\gamma - a, = \beta - \beta'$, also positive; and hence, the entropies on the plane of $xy$ are real; the umbilicus centres are imaginary for this plane, but real for the plane of $yz$ as appears next being
\[
a^2 x^2 = -\tfrac{a^2}{-\beta\gamma}b^3,\quad\text{Coordinates of Umbilicus centres in plane of $xy$ (real).}
\]
$b^2 y^2 = $ \ldots

\clearpage

% --- p.329 ---
\bigskip
\textit{p.~329 [520]}
\smallskip

\subsubsection*{Nodes of the Evolute.}

\textbf{23.} The Evolute is a curve with four nodes, all of them imaginary; viz.\ for the evolute in the plane of $xy$, the equation of which is
\[
(a x^2)^2 + (b y^2)^2 - c^2(a^2 y^2 + b^2 x^2) y^2 = 0,
\]
these are
\[
a^2 x^2 = -y^3,\quad b^2 y^2 = -a^3,\text{ Coordinates of Nodes of evolute in plane of $xy$ (imaginary)},
\]
in fact these values satisfy as well the equation of the evolute, as the two derived equations
\[
6a^2 x^2 [c^2(a^2 + b^2 y^2)^2] = b c y^2(a^2 + b^2 y^2) = 0,
\]
\[
6b^2 y^2 [c^2(a^2 + b^2 y^2)^2] = a c x^2(a^2 + b^2 y^2) = 0;
\]
or its points in question are nodes of the evolute.

The evolute has the four cusps on the axes and two cusps at infinity, in all 8 cusps as just remarked; it has 4 nodes; and the order being 6, the class is
\[
30 - 2 \cdot 4 = 8 \cdot 6 = 4.
\]

\subsubsection*{Section by the plane infinity.}

\textbf{24.} The surface itself is finite, and the section by the plane infinity is therefore imaginary; but by what precedes the nodal curve has nodal points at infinity, viz.\ there must be real acnodal points on this imaginary section. The section by the plane infinity resembles in fact the principal sections; viz.\ writing successively $\xi = \infty$, and $\eta = \infty$, we have
\begin{align*}
-\beta\gamma x^2 &= a\eta^4(\eta + a), & -\beta\gamma x^2 &= a\xi^4(\xi + a),\\
-\gamma a y^2 &= b\eta^4(\eta + b), & \text{or}\quad -\gamma a y^2 &= b\xi^4(\xi + b),\\
-a\beta z^2 &= c\eta^4(\eta + c), & -a\beta z^2 &= c\xi^4(\xi + c);
\end{align*}
giving respectively
\[
a^2 x^2 + b^2 y^2 + c^2 z^2 = \eta^4,\quad (a x^2)^2 + (b y^2)^2 + (c z^2)^2 = \xi^4,
\]
where the first equation represents an imaginary conic which counts three times; and the second equation, the rationalised form of which is
\[
(a^2 b^2 c^2)^2 + (b^2 c^2 a^2)^2 + (c^2 a^2 b^2)^2 - 2a^2 b^2 c^2 (a^2 y^2 + b^2 z^2 + c^2 x^2) z^2 + 8a b c \cdot a b c \cdot a b c \cdot y^2 z^2 = 0,
\]
is an imaginary evolute. The conic and evolute have four contacts and four simple intersections (in all $4, 2 + 4 = 12$ intersections) which are all of them imaginary. But the evolute has four real nodes (analogue of $a x^2 = -y^3$, etc.), which are repeating but being also by what precedes, the tangent of the evolute at that point, which point is also a node on the surface.

\clearpage

% --- p.330 ---
\bigskip
\textit{p.~330 [520]}
\smallskip

$a^2, y^2, z^2$, the lines in question will be not surely parallel to, but will be, the asymptotes of the nodal curve.

The plane infinity may be reckoned as a principal plane, and we may regard it in each of the four principal planes there are four umbilicar centres, four entropies, and four evolute-nodes.

\subsubsection*{The generation of the surface considered geometrically.}

\textbf{25.} I have deferred until late the discussion of the generation of the centro-surface by means of the centro-curves, for the reason that it can be carried on more precisely now that we know the forms of the principal sections and of the nodal curve. The two figures exhibit (as regards one extent of the surface) the portions already distinguished as (A), and (B): they intersect each other in the nodal curve, shown in each of the figures.

\begin{center}
\textit{[Figure: portions (A) and (B) diagram --- omitted]}
\end{center}

\textbf{26.} Consider first the generation of the portion (A) by means of the major-mean sequential centro-curves. The major-mean curves of curvature (extending to those below the plane of $xy$) commence with a portion (extending from umbilicus to umbilicus) of the ellipse $\tfrac{X^2}{a} + \tfrac{Y^2}{b} = 1$, this may be termed the vertical curve, and they end with the whole ellipse $\tfrac{X^2}{a} + \tfrac{Y^2}{b} = 1$, this may be termed the horizontal curve. The normals at the several points of the vertical curve successively intersect each other so as to form a circle, namely a part of the cuspidal curve, viz.\ from cusp $(b) (b)$ all the way to other umbilicar centres at once, viz.\ to a first-cusped curve in fig. $(A)$, which may be defined as the principal entropy of $-\beta\gamma$, the corresponding entrolicus being the the umbilicar centre $-\gamma a$. The successive forms of the major-mean entropy as we pass to the plane infinity are entropies as shown in the figure $(B)$, We have in each case attended only to the curve of curvature lying above the plane of $xy$ and the corresponding major-mean separated centro-curve is now in $\eta = a$; for to value of $\beta$ between $-c, -b$ inclusive; the resulting curve of curvature is for $\eta = -b$, the umbilicar shown above in fig. $(A)$, when these are only $\beta = b$, $\eta = a$, the umbilicar centre at first.

\clearpage

% --- p.331 ---
\bigskip
\textit{p.~331 [520]}
\smallskip

opening out the two component parts thereof: the two upper cusps of the curve $(b)$ are atomic on the geodesic of the curve-surface, and the two lower cusps upon two detached portions respectively of the cuspline of the centro-surface. And as the curve

\begin{center}
\textit{[Figure: A, B, C centro-curves diagram --- omitted]}
\end{center}

of curvature gradually broadens out and ultimately coincides with the $XY$-section of the ellipsoid, the four cusped curve continues to open itself out, and ultimately coincides as shown figure $(c)$ with the $xy$-evolute of the centro-surface, viz.\ this evolute is the sequential centro-curve belonging to the horizontal curve of curvature. The curves of curvature so that we have, for the special case which is now under consideration, the normal at any point of any of the major-mean separated centro-curve as the line tangent thereto. Each, as we have above said, of these normals meets in some entropy in the figure $(c)$.

\textbf{27.} We consider next the generation of the portion (B) by means of the major-mean concomitant centro-curves. Starting as before with the vertical curve of curvature, the concomitant centro-curve is a finite portion (terminated each way at an umbilicar centre) of the curve of the curve $XY$-evolute, viz.\ we have then the entropies $\xi, \xi$, which terminate the curve. The successive forms of the concomitant centro-curve are as shown in figure (B), We have in each case attended only to the curve of curvature below the plane of $xy$ and the corresponding sequential centro-curve is also more on the negative side. Pass we further to the cuspidal $X, Y, Z$ and lastly $(X, Y, Z) = $ the cuspidal $X, Y, Z$ as $X, Y, Z, X, Y, Z$, will completely deduced, the geometrical signification is anything but clear.

\textbf{28.} There is a precisely similar generation of the portion (B) by means of the minor-mean sequential centro-curves, and of the portion (B) by means of the minor-mean concomitant centro-curves.

\clearpage

% --- p.332 ---
\bigskip
\textit{p.~332 [520]}
\smallskip

\subsubsection*{The Nodal Curve. Art.\ Nos.\ 29 to 60.}

\textbf{29.} As before different points on the ellipsoid correspond to the same point on the centro-surface, this will be a point on the Nodal Curve: the conditions for this if $(\xi,\eta)$, $(\xi_0, \eta_0)$ are the parameters for the two points on the ellipsoid, are obviously
\[
a^2(a + \xi)^3(a + \eta) = a^2(a + \xi_0)^3(a + \eta_0),
\]
\[
b^2(b + \xi)^3(b + \eta) = b^2(b + \xi_0)^3(b + \eta_0),
\]
\[
c^2(c + \xi)^3(c + \eta) = c^2(c + \xi_0)^3(c + \eta_0);
\]
these equations in effect determine $\eta$ as a function of $\xi$, so that the three equations will be in effect a known function of $\xi$.

The relation between $\xi$ and $\eta$ in fact be found by eliminating $\xi_0, \eta_0$ from the foregoing equations, but it is easier to eliminate $\eta$ and $\eta_0$; thus consider obtaining $\xi$, and $\xi_0$ between as values of $\eta$, viz.\ may be regarded as a known function of $\xi$; in the equation in two equations may be expressed in terms of $\xi$, $\xi_0$, so that each of these quantities will be in effect a known function of $\xi$.

\textbf{30.} The relation between $\xi, \xi_0$ in the first instance given in the form
\[
0 = \begin{vmatrix} a^3(a + \xi)^3 + b^3 + c^3, & (a + \xi_0)^3, & (a + \xi)^3 \\
b^3(b + \xi)^3 + a^2 + c^2, & (b + \xi_0)^3, & (b + \xi)^3 \\
c^3(c + \xi)^3 + a + b, & (c + \xi_0)^3, & (c + \xi)^3 \end{vmatrix} = 0.
\]

Throwing out a factor $(\xi - \xi_0)$, this becomes
\[
\Sigma\{a^2(a c + b^2)^2 + a^2 + b^2 + c^2\}
\]
\[
+ (b - c)\{abc \, c^2 + (b a c)^2 + c^2(b + c)(b + c)(a + \xi)\} = 0,
\]
where the left-hand side is a symmetrical function of $\xi$ and $\xi_0$, and therefore divisible by $(\xi - \xi_0)^2$. It is obviously to $(b - c)$, viz.\ we have thus for the determination of $\eta$, $\eta_0$, the simple equations, obtained by throwing out the factor $\xi - \xi_0$.

To work this out, write $\xi + \xi_0 = p$, $\xi\xi_0 = q$, the result the rentals quotient being $\frac{1}{\Delta(p^2 - 4q)}$,
\[
\Sigma[(b - c)(a + p) + bc] = \begin{vmatrix} 3p^2 c, & 3 (a^2)(b^2 + c)(c^2)\\
b + cp & b\\
+ p^2(b + c) & c^2\\
+ b^3(b + c) & b\\
+ 3p^4 q & c \end{vmatrix} = 0,
\]
where the left-hand side divides by $\tfrac{1}{\Delta}(p^2 - 4q)$.

\footnotesize
$^1$ This was my first method of solution; and I have thought the results quite interesting simply to ensure them; but it will appear in time that I have succeeded in expressing $\xi, \xi_0, \eta, \eta_0$ in terms of a single parameter $\xi$.

\clearpage

% --- p.333 ---
\bigskip
\textit{p.~333 [520]}
\smallskip

\textbf{31.} Developing and reducing, and omitting this factor, the final result is
\[
6R + 3Q(p + P)(p + Pp) + 3pq = 0,
\]
when as before $P$, $Q$, $R$ denote $a + b + c$, $bc + ca + ab$, $abc$ respectively; that is,
\[
6R + 3Q(P + Pp) + P(P^2 + 4Q)q + 2pq + (P + Q)q = 0,
\]
or, as this may also be written,
\[
6R + 3Q(p + P_0) + (P + Q)q
\]
\[
+ \xi_0 (P + 4Q)
\]
\[
+ \xi_0^2 (P + Q) - 2\xi_0 q = 0,
\]
viz.\ the parameters $\xi, \xi_0$ have a symmetrical $(2, 2)$ correspondence.

\textbf{32.} From the equations $(a + \xi)^3 (a + \eta) = (a + \xi_0)^3 (a + \eta_0)$, etc., we have
\[
\Sigma (b - c)\{(a + \xi)^3 (a + \eta) - (a + \xi_0)^3 (a + \eta_0)\} = 0,
\]
and observing that the term in $(\eta - \eta_0)$ is
\[
\sqrt{[\beta(\beta - \gamma)(a + \xi)^2 + a(b - \gamma)(a + \xi)^2 + a(b - \gamma)(b + \xi)^2 + c(a - \beta)(c + \xi)^2]}
\]
\[
+ a(\gamma - \beta)\{2a + \xi\} - 2a\{a^2 + \xi\}
\]
\[
+ b(a - \gamma)\{2b + \xi\}
\]
\[
+ c(\beta - a)\{2c + \xi\},
\]
these are readily reduced to
\[
[3\xi - \xi_0 + (P + Q)p + (P + Q)(\xi + \xi_0)] \cdot \eta = -[(P + Q)\xi + 3\xi_0 + (b + a)\xi_0],
\]
\[
[(\xi - \xi_0) + a^2 P_0]\cdot (\eta - \eta_0) - \cdots,
\]
or, what is the same thing,
\[
[3(\xi - \xi_0) + 2P(\xi + \xi_0) + (P + Q)q + (P + Q)(\xi + \xi_0) + Q(\xi - \xi_0)] \cdot \eta = -\cdots,
\]
\[
3(\xi - \xi_0) + 2P(\xi + \xi_0) + Q(\xi - \xi_0) = -[(P + Q)\xi + 3\xi_0 + (b + a)\xi_0],
\]
and if we hence determine the values $3(\xi^2 - \xi_0^2) = \eta_0$, we find that the resulting terms divide by $\xi - \xi_0$, and we have
\[
3\xi_0^2(\xi_0 - \xi) - \xi^2[(b + a)\xi + 3\xi_0] = (\xi + \xi_0)\xi_0(b + a)(\xi + \xi_0)(\xi + a)
\]
\[
+ \xi_0 \cdots,
\]
\[
+ \xi_0 \cdots.
\]

Hence observing that by the relation between $\xi, \xi_0$, the first term is
\[
= -3Q\xi_0 + (P + Q)q + 2\xi_0 q,
\]
the equations become
\[
1: \eta : \eta_0 = -[Q + p(P + Q) - 3\xi_0] \cdot Q
\]
\[
: R(2\xi_0 + 2) - \xi^2(P + Q\xi)
\]
\[
: R(3\xi_0^2 - \xi^2) - \xi^2 P(P + Q\xi).
\]

\clearpage

% --- p.334 ---
\bigskip
\textit{p.~334 [520]}
\smallskip

and we thus have
\[
\eta = \frac{R(\xi_0^2 - \xi^2) - \xi^2 P(P + Q\xi)}{P\{Q + q(P + Q) - 3\xi_0\}},
\]
which, considering $\xi_0$ as a given function of $\xi$, gives $\eta$ as a function of $\xi$.

\textbf{33.} I write $\xi - \xi_0 = 2u$, $\xi + \xi_0 = 2t$, so that $p = 2t$, $q = t^2 - u^2$, $\xi = t + u$, $\xi_0 = t - u$. The relation between $\xi, \xi_0$ thus becomes
\[
6R + 6Qt + (P + Q)(t^2 - u^2) = 0,
\]
or, what is the same thing,
\[
u^2 = t^2 + \tfrac{6R + 6Qt}{P + Q};
\]
so that taking $t$ as a parameter, and considering $u$ as a node real value of $u$, respectively, we have the relation
\[
\xi = \tfrac{(t + u)\sqrt{P + Q} + R(P + Q\xi + Q\xi_0)}{P + Q + 2t},
\]
and the value of $\eta$ is taken as $-3\xi, \xi_0, -2 = -2\xi$. The relation between $\xi, \xi_0$ in the simple form $\xi + \xi_0 = 2t$, hence appears as an arithmetic mode of expression: which is in some respects advantageous, but I propose that this can be modified.

\textbf{34.} On comparison of reality is easily found
\[
y^2 = \frac{(a + a^2)(a + b)\{a + bc(a + 1)\} + a c(a + 1)\{a + b - c\}}{a + b + c\{a + (b + c)\}},
\]
where $a = -\tfrac{1}{2}(P + a)$, $b = \tfrac{1}{2}(P + Q)$; viz.\ as we then have $a$ in entrunctate coordinates of a point in a plane, $(a, t)$ will be the rectangular coordinates of the same point referred to as $a$, and also the equation of conic referred to set $(2x - 2, x)$ at $(a, t)$, viz.\ we have for the entire centro-surface in this way.

\textbf{35.} The curve is a conic curve symmetrical in regard to the axis of $x$, and having the three asymptotes,
\[
x = -t,\quad y = \tfrac{1}{2}(P + a) \pm \tfrac{1}{2}(P + a)\sqrt{c(P + 1)(P - a)};
\]
moreover we have $y = 0$ for the values $x = a$, $x$, $C$ where the curve cuts the axes of $C$, $E$, $F$, as shown in the figure: $a$, $C$, $E$, are at the distances from the centre, and with these data it is easy to draw the curve. It is then the form $\tau$, that whence writing $t = \tau, u = 1$ to $\eta + a$, and let $t \cdots \tau + a$, corresponding to the umbilicar centre; and at $\Omega, \Omega$, we have $\xi = -\beta$, $\eta = a$, and $\xi - \beta$, etc. corresponding to the corner.

\clearpage

% --- p.335 ---
\bigskip
\textit{p.~335 [520]}
\smallskip

We have
\[
A^2 + B^2 c = \tfrac{1}{2}\Delta + a^2(P + 1);
\]
\[
B = -\tfrac{1}{2}(\gamma + a)^2(\eta + a) + 2\Delta(P + 2 + \gamma),
\]
\[
B = -\tfrac{1}{2}(\gamma + a)^2(b + a^2 \omega),
\]
hence the rational part in question is
\[
= \pm \tfrac{(\eta + a)^2}{(2a + Pp + a)}\sqrt{\tfrac{a + R + (1 - \xi)(2 + \xi)}{P\{Q + q(P + Q) - 3\xi_0\}}}(2 + \xi)
\]

\textbf{32.} Write $\xi - \xi_0 = a + \tfrac{a}{u}$, etc., $\xi - \xi_0 = -\tfrac{a}{2}, \xi_0 + 2 = \tfrac{2}{u}, c$, $y = 0; p = -a, c$. The relation between $\xi, \xi_0$ thus becomes
\[
\xi - \xi_0 = a + \tfrac{a}{2u}, \quad \xi - \xi_0 = -\tfrac{a}{2}, \xi_0 + 2 = \tfrac{2}{u}, c, y = 0; p = -a, c.
\]
or as that the curve $\Omega = 0$ is obtained by eliminating $\xi$ from this equation and the foregoing, regard to $\xi$ at the cubic function becomes
\[
(a^2 + \xi)(b - \xi)\!\left(c + \xi + \tfrac{\xi}{a + \beta}\right) = 0;
\]
that is, the equation has the values
\[
\xi = -a^2, \quad \xi = -\beta^2, \quad \xi = -\xi_0 = -a(a + \beta).
\]

But these last values give
\[
a + \xi = \tfrac{a a}{a + \beta}, \quad \beta + \xi = -\tfrac{a a}{a + \beta}.
\]

\clearpage

% --- p.336 ---
\bigskip
\textit{p.~336 [520]}
\smallskip

and for $\xi = ($ we have
\[
u(2(a + \xi) - 1)\!\left(\tfrac{2u\beta\gamma}{a + \beta}\right)\eta_0 = \tfrac{(a^2)}{(a - \beta)^3}\eta, \quad \tfrac{(a - \beta)^3}{a P\beta}\eta_0,
\]
so that solving for greater simplicity, $(a - \beta)\eta_0 = -a\beta\eta$, the formulae become
\[
\tfrac{2P\xi}{a^2} = \tfrac{3a^2}{a^2} = \tfrac{3a^2}{\beta - a},
\]
\[
\tfrac{2P\xi}{\beta^2} = \tfrac{3\beta^2}{\beta - a},
\]
\[
\tfrac{2P\xi}{\gamma^2} = \tfrac{-a\beta\gamma}{a^2 + a\beta} = -\tfrac{a + \beta}{a},
\]
\[
\tfrac{a + \beta}{\Omega \cdot 1}\Omega = \cdots,
\]

\textbf{32.} This shows that there is the curtsey a cusp, the cuspidal tangent being in the plane of $xy$. It appears moreover that this tangent constants with the tangent of the evolute. In fact, from the equation $(a x^2)^2 + (b y^2)^2 - a^2 y^2 = 0$ of the evolute we have
\[
\tfrac{(a x)^2}{a} + \tfrac{(b y)^2}{b} + \tfrac{(a x^2)^2 + (b y^2)^2}{a^2} = 0,
\]
or substituting for $(a, y)$ their values as in the entropy,
\[
\tfrac{2(a - a)\{b - a\}\beta - a}{x^2(a - \beta)^2} = \tfrac{a + a x}{\beta + a^2 - a^2}\cdot\tfrac{dy}{dx} = 0,
\]
that is,
\[
\tfrac{d}{dx} = \tfrac{a x^2 + a x^2(\beta - a)}{a^2(a + \beta) - a^2\cdot\tfrac{dy}{dx}}.
\]
which is satisfied by the foregoing values of $\tfrac{x}{a}$ and $\tfrac{y}{a}$ at the two tangents there coincide.

We have
\[
4(da^2 + (b a^2)(b - 2a^2 - \cdots\sqrt{4}\Sigma a y^2 + a x^2) = \tfrac{a + a + a}{a(a + \beta)^2 a}\eta\cdots\xi,
\]
which is virtue of
\[
(a y)(2 + y)(a + \beta) = \beta y^2(2y + a) + \beta y^2(2y + a) + (2y(2 + a) + a) - 4y^3(2 + a) - 3y^4 y^2(a + \beta),
\]
is
\[
4[(b y)^2 + (b y)^2 + (a y^2 + 3y^4 y^2(a + \beta))] = \tfrac{2(a + \beta)^2}{a(a + \beta)^2}\eta\cdots\xi.
\]
(observe $2(a^2 + \beta) - a(a - \beta), = -(a + (a + \beta), = -8a^2$, is negative)
\[
= -\tfrac{2a^2}{a(a + \beta)^2}\eta\cdots\xi.
\]

\clearpage

% --- p.337 ---
\bigskip
\textit{p.~337 [520]}
\smallskip

if $\tfrac{dy}{dx}$ be the value at the entropy. Writing $\tfrac{a^2}{a^2}$ for the element of the are we have
\[
\frac{ds^2}{a^2}\frac{dx^2(a - \beta)^3}{(a - \beta)^2} - \tfrac{4}{a}\eta - \tfrac{8a^2}{a},
\]
\[
\frac{ds^2}{a^2}\frac{dx^2(- a\beta\gamma)^2}{(a - \beta)^2} - \tfrac{8a^2}{a},
\]
which exhibit the form at the entropy.

\subsubsection*{The Nodal Curve, expressions for the coordinates in terms of a single parameter $a$. Art.\ Nos.\ 33 to 41.}

\textbf{33.} After the foregoing investigation of the nodal curve, I was led to perceive that it is possible to express $\xi, \eta, \xi_0, \eta_0$ in terms of a single variable $\eta$; and thus to obtain expressions for the coordinates of a point of the nodal curve in terms of the single variable $a$. The analytical form of the nodal curve being real, with its imaginary values of $\xi, \eta$, the modes $y$ may we may assume
\[
\xi = -b^2 - p(b^2 - a^2 - b^2),
\]
\[
\eta = -b^2 + p(b^2 - a^2 - b^2),
\]
$(i = \sqrt{-1}$ as usual): viz.\ $p$ denote one or other of the quantities
\[
y_+ = a + (a^2 - b^2),\quad q = a - (a^2 - b^2);
\]
the expressions for $-\beta\gamma x^2$, $-\gamma a y^2$ become
\[
= [(b - p)(a + 1)(b - q)^3](b - p)^3,
\]
and we have therefore the condition that this shall be real (for the two values $b = p$, $\dot{b} = -b$); being real, it will in certain cases be positive, in other cases negative, and the values for the remaining condition is given by the equation
\[
\Delta'(b^2 - \xi^2 - \xi_0^2 - 4)(b^2 + b^2 + \xi^2 + 2b^2)^2 + p^2(b - \xi^2 - 2)(b^2 + \xi)^2 - 0 = 0.
\]

\textbf{34.} The condition of reality is easily found to be
\[
(\beta - p)(p + \beta - 1) + (a + p)^2 + (\gamma + p)^2 + 4(a + p)(\beta - p) + (a + p)^2(b - \beta) - p^2 q^2 = 0;
\]
viz.\ this equation in $\Delta$ must have the form $-a, b$, suppose to a value
\[
P = p, \quad q = qP_0p + p^2,
\]
we have therefore
\[
\tfrac{(\gamma - a)\beta}{a^2} = a + \tfrac{a(a^2 - b^2)}{2bg(\gamma - a)}\sqrt{Y}\cdot\eta(a - \beta - \gamma)(a + \beta + \gamma),
\]
where, regarding $X$, $Y$ are given functions of $a$, the coordinates of a plane in a cubic curve having the point $X = 0, Y = 0$ for a node, writing $Y = b + 3a(X + 1)$, $X$ will be such of these is rational function of $a$. The curve has the property that being any line through the entropy, the tangents at the points of the curve, the corresponding entropies on it are
\[
(a^2 + Y)(a^2 + a)(a + 1) + (b + 1)(a^2 + b)\beta(a^2 + a) = -a^2(b - p)^3 + a(p + p)^3,
\]
for a point on the nodal curve.

\textbf{35.} To complete the investigation, writing for a moment $T = 3 = 3a^2(X + 1)$, we obtain
\[
\frac{(b - p)^2}{a^2} = \tfrac{a}{2a^2(a + a - 1)} + \tfrac{a(a^2)(b - \xi)^3}{a + b - 1},
\]
or putting for a moment
\[
\frac{(\gamma - a)\beta}{a^2} = K,
\]

\clearpage

% --- p.338 ---
\bigskip
\textit{p.~338 [520]}
\smallskip

we have
\[
X = \tfrac{a(K - 7)(a - 1)}{a + 1 - K},
\]
\[
X_0 = \tfrac{K(a - 7)(a + 1)}{a + 1 - K},
\]
\[
Y = K \tfrac{3K(a^2 - 1) + 2(a + 1)}{a + 1 - K},
\]
\[
Y_0 = K \tfrac{3(a - 1)(K - 2)(a + 1)}{a + 1 - K},
\]
\[
Z = 2K \tfrac{(b - 1)X(a - 2 + 1)}{a + 1 - K},
\]
\[
Z_0 = 2K \tfrac{(b - 1)(X - 2 + 1)}{a + 1 - K},
\]
or substituting for $K$ its value we have
\[
K(a - 1) = \tfrac{(a - 1)^2}{ab}\!\left(a^2 - \tfrac{8aq}{(a - 1)^2}\right),
\]
\[
= -\tfrac{(a - 1)^2}{ab}\!\left(a + \tfrac{2q}{a - 1}\right)\!\left(a - \tfrac{2q}{a - 1}\right);
\]
$a a + 1 - K = -\tfrac{1}{ab}(a^2 + 1)\!\left(a + b - 1\right)(a - p), \cdots$
$+ b\!(a - p)$,
and changing the sign of the radical we have the values of $\tfrac{x}{a}, \cdots$,

\textbf{36.} Write for a moment,
\[
\left[\Delta - \tfrac{1}{2}(\gamma - a)\sqrt{1 + \sqrt{\tfrac{2a(a - 2 - 4\beta + 2)}{a(a - 2)}}}\right] - (a - 8a\beta + \sqrt{\beta})Y + B\sqrt{\delta} = 0,
\]
\[
\left[\Delta - \tfrac{1}{2}(\gamma - a)\sqrt{1 + \sqrt{\tfrac{2a(a - 2 - 4\beta + 2)}{a(a - 2)}}}\right] - (a - 8a\beta + \sqrt{\beta})Y + B\sqrt{\delta} = 0,
\]
then in the product of these two expressions the rational part is $= A^2A + B^2B$; and from the manner in which they were arrived at we have $9 = A_1A + B_1B$, and the rational part is thus
\[
= -\tfrac{16}{a^2}(B^2 - B'B).
\]

\clearpage

% --- p.339 ---
\bigskip
\textit{p.~339 [520]}
\smallskip

% TODO: book page 339 (very dense algebraic manipulation; transcription unreliable)

We have
\[
A^2A - B^2B = (\Delta - a)^2 X - 8a^2 X_0,
\]
\[
B' = \tfrac{1}{2}(\gamma + a)X\!\left[a - \tfrac{2aq}{(a - 1)^2}(X + S)\right] - a^2 b^2(P^2 - 2a^2)
\]
\[
- B(\gamma + R)\!\left[(\gamma - a)(b - \beta) + Y(\gamma - a) - 2y - \tfrac{4ay}{a^2}\right],
\]
or, what is the same thing,
\[
\Delta\beta + \xi\!\left(1 + \tfrac{q}{a - 1}\right) = \xi a P + 2(P - a)X.
\]

The first thus becomes
\[
\pm \xi(a + bX)\!\left(\tfrac{a + b X}{(a + b - 1)\beta - a^2}\right) - X(a + b)(b - a)X - 4a + 4(b + 1)
\]
which is
\[
= -aX - X(P - a) - (b + 1)(a^2X^2 - 2a + 4(b + 1)X);
\]
and for the value of $\eta$, proceeding to the third term, this is
\[
= -aX_0 + X(P - a) - (b + 1)(a^2 X_0^2 - 2X_0 + 4(b + 1)X_0);
\]
so that without any further reduction, $\eta_0 = -\beta\eta$.

\textbf{37.} We have
\[
\tfrac{-2bx}{x^2}(a - 1) - \tfrac{1}{a(a - 1) - 1}\!\!\left[3a a^2 - 2\beta - \xi\!\left(a^2 - \tfrac{4qa}{(a^2 - 1)^2}\right)\right]
\]
\[
- \tfrac{k}{(a^2 - 1)(a - 1)}\!\left[3a a^2 - 2\beta - \xi - \tfrac{4qa}{(a^2 - 1)^2}\right] + \cdots,
\]
where the term in [\ ] is
\[
- \tfrac{2bg}{a^2(a^2 - 1)}\!\left[3a a^2 - 2\beta - \xi\!\left(a^2 - \tfrac{4qa}{(a^2 - 1)^2}\right)\right],
\]
$3\eta + 2 + (b a^2 \Sigma X)(a^2 + 1) - 2\beta + 4 = - ab(a^2 - 1)\xi - 2bq(a^2 - 1)y - q^2$
$+ (q a^2)(a^2 - 4)$,
where in the case the first term is
\[
= 3a a^2(a + 1)(a^2 + 1) + 4(\sigma + \beta)(b - \beta), \quad \cdots ,
\]
\[
+ 8(\beta + a)(\beta - a),
\]
\[
+ (b + 1)(b + 1)(a - 1)^2 - q.
\]

Hence observing that $\Delta$ in like manner multiplying by $(\beta - q)^2$, we have
\[
-\tfrac{2bx}{x^2}(a - 1) = -\tfrac{a(a + b)(a^2)(3a - 2\beta - 4q\beta + 8\beta(a - q + 1) + 8a(a^2 - 1))}{a^2(a - 1)(a^2 - 1)} - \tfrac{a^2(b - a)}{(a^2 - 1)}\!\left[3a a^2 + 4(b - \beta) + 8(\beta - q) + 8(a^2 - 1)\right] + \tfrac{a(b + a)}{a^2(a - 1)}\Sigma,
\]

\clearpage

% --- p.340 ---
\bigskip
\textit{p.~340 [520]}
\smallskip

% TODO: book page 340 (dense algebraic manipulations)
\[
[-2at(\beta - a) + 5\beta\beta(\beta + 2)\beta\beta + 4(a - 1)\beta\beta + \Omega - 3\beta\beta(a^2 - 1)\beta\beta + (\beta - 1)\beta\Omega] - \tfrac{2\beta y}{(a^2 - 1)}(a^2 - p)(a^2 - q),
\]
or finally
\[
\Omega : \beta(p - q) + 2(a - 1)(\Omega) = (- 2\beta(P + 1) - \beta - q)(\beta - p)(\beta - q).
\]

But $a^2 + 4a, \xi + \xi_0 = \xi$, viz.\ since $\xi - \xi_0 + a^2(2a + 1)$, and therefore
\[
a^2(P - 1) + 2 + 8 - \beta = 2a^2(\xi^2 + \xi_0)(\xi + \xi_0),
\]
the equation thus divides by $-2a^2(P + 1) + a^2 - \xi - \beta + 2y - a^2$. So that $\Omega$ has the value originally so denoted, we have
\[
\eta = -a^2 + \tfrac{8\beta\beta}{(a - 1)^2}\beta,
\]

\textbf{38.} Lastly the equation $\Omega = \xi'(b + \xi)\!\left(c + \xi + \tfrac{\xi b\beta\gamma}{a + \beta}\right) = 0$ is satisfied if $a^2 + \beta\xi = 0, q + b - \beta$, the equations
\[
\Omega + \xi'(b + \xi)\!\left(c + \xi + \tfrac{\xi b\beta\gamma}{a + \beta}\right) = 0,
\]
\[
(a^2 + \xi)(b + \xi) + a^2(b + \xi) + \xi b\beta\xi = 0,
\]
then give
\[
(a^2 + \xi)(b + \xi)(c + \xi) = \xi a\beta\gamma,
\]
which can be satisfied by $\xi = a$, giving $\eta = a + \beta = -a^2$, which in the case I, or else by
\[
\xi'(b + \xi) - \xi^2 - \beta a\beta\gamma,
\]
\[
\xi'(b + \xi) - \xi^2 + \xi\beta a = a^2,
\]
that is,
\[
\xi(a + \beta) - \xi^2 - a^2,
\]
\[
\xi(b - \beta) + \beta^2 - a^2.
\]

To show that these values satisfy the relation between $\xi, \xi$, observe that they give
\[
\xi + \xi_0 = a + b - 2a^2,\quad \xi\xi_0 = b - \beta - a^2,
\]
whence also
\[
\xi'\xi_0 = \xi b - \xi^2 - a^2,\quad \xi - \xi_0 = \xi(a - \beta + 1) - 1,
\]
and the relation becomes
\[
6a\beta\gamma - 3(b + ac + a^2 + a^2)(a + b - \xi^2) + (a + b + 1)(a^2 + b - \xi^2 - 1) = 0,
\]
which is an identity.

\clearpage

% --- p.341 ---
\bigskip
\textit{p.~341 [520]}
\smallskip

\textbf{43.} I will show that these values of $\xi, \xi_0$ give the foregoing values $\eta = \eta_0 = -a^2$. We have
\begin{align*}
1 : \eta - \eta_0 &= P\{Q, q(P + Q) - 3\xi_0\} : 5R \\
&: (\xi_0 - \xi) [3R - (P + Q)(2\xi_0 - q) + \xi P\beta_0 + (2\xi_0 - q)] \\
&: (\xi_0 - \xi) [3R - (P + Q)(2\xi_0 - q) + \xi^2 P + (P + 2\xi_0 - q)];
\end{align*}
this is
\[
1 : \eta - \eta_0 = P\{Q, q(P + Q) - 3\xi_0\} : 5R\!\left(\xi P - \tfrac{2}{a}\right) + \xi^2 P + \xi(P - \xi)
\]
or
\[
\eta - \eta_0 = -a^2,\quad \eta + \eta_0 = -2a^2,\text{ that is, } \eta = \eta_0 = -a^2.
\]

\textbf{44.} For the real umbilicar centres and entropies we have

I. $\xi = -b^2,\quad \eta = a^2,\quad \xi_0 = -b^2,\quad \eta_0 = -b^2;\text{ we have}$
\begin{align*}
X^2 &= a^2(a^2 - b^2)/a, &
X_0^2 &= -a^2/a, \\
Y &= 0, & Y_0 &= 0,\\
B &= 0, & E_0 &= -c^2/c,\\
a^2 X^2 &= \tfrac{(a - b)^2}{a^2},\\
y &= 0, & \text{(real umbilicar centre)},\\
c^2 z^2 &= \tfrac{a^2 b^2}{a^2}.
\end{align*}

II. $\xi = -b^2, \quad \eta = a(a + \beta) = -2b^2$:
\[
(a^2 + \eta = -a\beta) - a\eta = b - \beta - \beta,\quad b + \eta = \tfrac{a(a - b)}{a^2}, \quad c + \eta = \tfrac{c(b - \beta)}{(a^2 - \beta)^2},
\]
\[
\xi_0 = -b^2,\quad \eta_0 = \tfrac{2b\beta}{a^2 - \beta}, \quad \xi_0 + \eta_0 = \tfrac{a(\xi - \beta)}{a^2 - \beta}, \tfrac{b - \beta}{a^2 - \beta}, \tfrac{c + \beta - \beta^2}{a^2 - \beta};
\]
\[
b + \eta_0 = -\tfrac{a^2}{(a - \beta)^2},
\]
\begin{align*}
X^2 &= \tfrac{a^2}{a^2 - \beta^2}\,\tfrac{(b - \beta)^3}{a - \beta}, &
X_0^2 &= -\tfrac{a^2}{a^2 - \beta^2}\,\tfrac{(b - \beta)^3}{a - \beta},\\
Y^2 &= -\tfrac{b^2}{a^2 b^2}\,\tfrac{(a - \beta)^3}{a^2 - \beta^2}, &
Y_0^2 &= \tfrac{b^2}{a^2 b^2}\,\tfrac{(a - \beta)^3}{a^2 - \beta^2},\\
Z &= 0, & Z_0 &= 0,
\end{align*}
ellipse concomitant,\hfill ellipse separated;
\begin{align*}
a^2 x^2 &= \tfrac{a^2}{(a - \beta)^3}\,\tfrac{(b - \beta)^3}{a^2 - \beta^2},\\
b^2 y^2 &= -\tfrac{b^2}{(a - \beta)^3}\,\tfrac{(a - \beta)^3}{a^2 - \beta^2}\quad\text{(real entropy)},\\
z &= 0.
\end{align*}

\clearpage

% --- p.342 ---
\bigskip
\textit{p.~342 [520]}
\smallskip

% TODO: book page 342 (highly intricate equations and identities; partial transcription)

\textit{Nodal curve in vicinity of umbilicus centre, $a x^2 = -y^2$, $c^2 z^2 = w^2$, Art.\ No.\ 47 to 49.}

\textbf{47.} Write
\[
\xi = -b^2 + p,\quad \eta = -b^2 + r,\quad \xi_0 = -b^2 + q, \quad \eta_0 = -b^2 + s;
\]
we have to find the relation between $p$, $q$, $r$, $s$; first for $p$, $q$, the equation of correspondence gives
\[
6R + 5Q(p + q + 2b^2 - 8b^2) + (P + Q)(p + q + b^2)
\]
\[
+ 5\{-2b^2 + 2bQp + p^2 - b^2(p + q) - 2bp - p\} = 0,
\]
that is,
\[
3(p + q)(3b^2 - 2b^2) + (b^2 - b^2)(p + q + 2bp) = 0,
\]
or writing these equations as follows
\[
3(p + q)(\gamma - \beta) + (3 \beta + 3\gamma - 2)\beta\eta = 0,
\]
viz.\ this is
\[
3p + 3q + \tfrac{8\beta\eta}{\gamma - \beta}(p + q) = 0;
\]
or in symbolic form
\[
3p + q + \tfrac{2\beta\eta}{\gamma - \beta}(\xi_0)^2 = 0.
\]

\textbf{48.} Now the equations
\[
(a^2 + \xi)(a + \xi)^2 = (a^2 + \xi_0)(a + \xi_0)^2 = a^2(b + \xi)^2(b + \xi_0)^2 = (a + \xi)^2(\cdots),
\]
putting therein for $\xi, \eta, \xi_0, \eta_0$ their values, give the first of them
\[
\log\!\left(1 + \tfrac{r}{a^2}\right) - \log\!\left(1 + \tfrac{s}{a^2}\right) = 8\log\!\left(1 + \tfrac{p}{a^2}\right) - 8\log\!\left(1 + \tfrac{q}{a^2}\right),
\]
that is,
\[
r + 3p + \tfrac{1}{2a^2}(r^2 + 3p^2) + \tfrac{1}{3a^4}(r^3 + 3p^3) + \cdots = +s + 3q + \tfrac{1}{2a^2}(s^2 + 3q^2) + \tfrac{1}{3a^4}(s^3 + 3q^3) + \cdots ;
\]
and similarly the second equation
\[
r + 3p + \tfrac{1}{2a^2}(r^2 + 3p^2) + \tfrac{1}{3a^4}(r^3 + 3p^3) + \cdots = s + 3q + \tfrac{1}{2a^2}(s^2 + 3q^2) + \tfrac{1}{3a^4}(s^3 + 3q^3) + \cdots;
\]

\clearpage

% --- p.343 ---
\bigskip
\textit{p.~343 [520]}
\smallskip

whence multiplying by $\gamma$, and adding,
\[
(\gamma + s)\!\left[r + 3p + \tfrac{1}{2}(r^2 + 3p^2)\right] = (\gamma + s)\!\left[r + 3p + 3p_0 + \tfrac{1}{3a^4}(r^3 + 3p^3)\right],
\]
which, neglecting terms of the third order, is
\[
r + 3p = s + 3q,
\]
Subtracting the two equations we have
\[
\tfrac{1}{3}\!\left(\tfrac{1}{a^2} + \tfrac{1}{b^2}\right)(r^2 + 3p^2) + \tfrac{1}{a^2}\!\left(\tfrac{1}{a^2} - \tfrac{1}{b^2}\right)(r^3 + 3p^3) = \tfrac{1}{3}\!\left(\tfrac{1}{a^2} + \tfrac{1}{b^2}\right)(s^2 + 3q^2) + \tfrac{1}{a^2}\!\left(\tfrac{1}{a^2} - \tfrac{1}{b^2}\right)(s^3 + 3q^3);
\]
viz.\ this is the same thing,
\[
r^2 + 3p^2 + \tfrac{2}{3}\!\left(\tfrac{\gamma - a^2}{\gamma}\right)\!(r^3 + 3p^3) = s^2 + 3q^2 + \tfrac{2}{3}(s^2 + 3q^2)(r - s) - 0,
\]
or, what is the same thing,
\[
r^2 - s^2 + 3(p^2 - q^2) + \tfrac{1}{3}\!\left(\tfrac{\gamma - a^2}{\gamma}\right)\!(r^2 + s^2)(r - s) - 0,
\]
which, putting therein $r - s = -3(p - q)$, is
\[
-r - s + 3(p + q) + 3a\!\left(\tfrac{\gamma - a^2}{\gamma}\right)\!(r^2 + s^2) + 9(p^2 + q^2) = 0,
\]
say this is
\[
-r - s + 3(p + q) + 2\Delta = 0,
\]
combining herewith
\[
r - s + 3(p - q) = 0,
\]
we have
\[
r + 3p - \Delta = 0,
\]
and
\[
s + 3q - \Delta = 0,
\]
where
\[
\Delta = \tfrac{1}{2}\tfrac{\gamma + a^2}{\gamma}(-r^2 - s^2 + p^2 + qq + q^2),
\]
But substituting herein the values $r + 3p = -3p$, $s = -3q$, this becomes
\[
\Delta = \tfrac{1}{2}\!\left(\tfrac{\gamma + a^2}{\gamma}\right)\!\!\left[-4p^2 - 4q^2 - 3p^2\right] = -\tfrac{1}{2}\!\left(\tfrac{\gamma + a^2}{\gamma}\right)\!q^2,
\]
and then
\[
r = -p + 3p_0 + \Delta = -3p - \tfrac{1}{2}\!\left(\tfrac{\gamma + a^2}{\gamma}\right)\!q^2,
\]
that is,
\[
r = 3p + 2(p + q) + \Delta = -\tfrac{4(\gamma - a)}{\gamma a^2}q.
\]

\clearpage

% --- p.344 ---
\bigskip
\textit{p.~344 [520]}
\smallskip

\textbf{49.} We have then
\[
a^2 x^2 = \tfrac{y^2}{a^2}\!\left(1 + \tfrac{r}{a^2}\right)\!\left(1 + \tfrac{p}{a^2}\right) = \tfrac{y^2}{a^2}\!\left(1 + \tfrac{r + 3p}{a^2} + \tfrac{3p(r + p)}{a^4}\right) = \tfrac{y^2}{a^2}\!\left(1 + \tfrac{-4(\gamma - a)q}{ay\gamma} + \tfrac{6q}{y a^2}\right) = \tfrac{y^2}{a^2}\!\left[1 + \tfrac{2(p - 3q)}{a\gamma y}\right],
\]
and in the same way from $a^2 x^2 = \tfrac{a^2}{b^2}\!\left(1 - \tfrac{r}{a^2}\right)\!\left(1 - \tfrac{p}{a^2}\right)$, we have
\[
a^2 x_0^2 = \tfrac{y^2}{a^2}\!\left[1 + \tfrac{2(p + \beta)}{a\gamma}\right];
\]
moreover we have at once
\[
b^2 y^2 = \tfrac{a^2}{a^2} - \tfrac{1}{a\gamma}q^2.
\]

Hence, writing $x + \delta x, 0 + \delta y, z + \delta z$ for $x, y, z$, we find
\[
\delta x = \tfrac{1}{2}p \tfrac{2(p - 3q)}{a y\gamma}\cdot q^2,
\]
\[
\delta y = \tfrac{1}{2}p\sqrt{\tfrac{a}{\gamma}}\cdot q,
\]
\[
\delta z = \tfrac{1}{2}p \tfrac{2(p - \beta)}{a y\gamma}\cdot q^2,
\]
or, what is the same thing,
\[
\delta x : \delta y : \delta z = \tfrac{2(p - 3q)}{a y\gamma} : \sqrt{\tfrac{a}{\gamma}} : \tfrac{2(p - \beta)}{a y\gamma},
\]
where $a, z$ denote the value at the umbilicar centre.

\subsubsection*{Nodal curve in vicinity of real entropy, viz.\ $a x^2 = \alpha^2, b y^2 = \beta^2$. Art.\ Nos.\ 50 to 52.}

\textbf{50.} Write
\[
\xi = -a^2 + p,\quad \eta = -a^2 + r,
\]
\[
\xi_0 = -a^2 + q,\quad \eta_0 = -a^2 + s.
\]
% TODO: book page 344 dense (entropy neighbourhood, page partially transcribed)

\clearpage

% --- p.345 ---
\bigskip
\textit{p.~345 [520]}
\smallskip

% TODO: book page 345 dense (continuation)
and first for the relation between $\eta$ and $\eta_0$, writing for a moment $\tfrac{3p\eta}{a - \beta} = \eta_0 = -Q$, and therefore $\xi_0 = -a + Q$, the equation of correspondence gives
\[
-3a\beta\gamma + (\eta_0 + s)\beta + (\beta + \eta + \tfrac{1}{2}\eta^2) - 3\eta\eta(\eta + Q) - 3a^2(s + Q) = 0,
\]
which, putting for $Q$ its value, is
\[
-9\eta(\eta - s + 3rs - \tfrac{6a^2}{a - \beta})
\]
\[
+ (s + 3\eta r)\!\left[a^2(s + r) - \beta(r + r) + (a + \beta)\eta s - \tfrac{3a\beta}{a^2 - \beta}(\eta + r)\right]
\]
\[
- 3a\beta(s + r) - 3a^2(s + r) - \tfrac{4a^2 \beta}{a - \beta} = 0;
\]
or, what is the same thing,
\[
(3\eta s + 2\eta^2(r - \beta)) - 4a^2\eta s
\]
\[
+ (s^2 + 7a\eta + 5\eta s) + \tfrac{a^2 - 4a\beta + \beta^2}{a - \beta} + 4(\eta + s)\eta a + s - \beta
\]
\[
+ 3a(\eta + s) = 0,
\]
or, in small values,
\[
\left(2a + \tfrac{a + \beta a^2}{a - \beta}\right)(s + 3r) = 2a\beta r.
\]
viz.\ writing $\tfrac{a}{(a - \beta)^2}$ for $\tfrac{a}{(a - \beta)^2}$, this becomes
\[
\tfrac{2a^2}{a - \beta}(s + 3r) = 2a\beta(\eta - s).
\]

\textbf{51.} Moreover, from the equation $(c + \xi)^3 (c + \eta) = (c + \xi_0)^2(c + \eta_0)$, we have $c$ since $a$ and $a$ are the same order: $\beta$, $a$ of the order $\epsilon$; arising from the variation of $a$, an indefinitely small in regard to those arising from the variation of $\xi$; and we have
\[
\tfrac{3\eta s}{a - \beta} = -\tfrac{3a\beta}{a^2}(s - r) = -3p\eta - \tfrac{(p - q)}{2(a - \gamma)},
\]
\[
\tfrac{3\eta s}{a - \beta} = -3a - \tfrac{3a\beta}{a^2}\eta(p - q) = -3p\eta - \tfrac{(p - q)}{2(a - \gamma)},
\]

\clearpage

% --- p.346 ---
\bigskip
\textit{p.~346 [520]}
\smallskip

and for $\delta x$ etc.\ we have
\[
a^2 x^2 = -\tfrac{1}{a^2}\!\left(\tfrac{a\beta\gamma}{a^2}\right)^2 p\eta = -\tfrac{(a^2 - \beta)^2}{(a - \beta)^2}\!\left[\tfrac{-(a - \beta)^2}{4(\gamma + 2)\eta\eta}\right],
\]
so that solving for greater simplicity, $(\beta - a)\eta_0 = -y\beta$, the formulae become
\[
\tfrac{2a^2}{a - \beta}(s + 3r) = -\tfrac{(b - a)^2 \, \delta\eta}{a - \beta} - \tfrac{(a - \beta) \, \delta\eta}{(a - \beta)^2 \, \beta} q^2.
\]

\textbf{52.} This shows that there is at the entropy a cusp, the cuspidal tangent being in the plane of $xy$. It appears moreover that this tangent coincides with the tangent of the evolute. In fact, from the equation $(ax^2)^2 + (by^2)^2 - c^2 y^2 = 0$ of the evolute we have
\[
\tfrac{(ax)^3}{a} + \tfrac{(by)^3}{b} - \tfrac{(ax^2)^2 + (by^2)^2}{a^2}\,\tfrac{dy}{dx} = 0,
\]
or substituting for $(x, y)$ their values as in the entropy,
\[
\tfrac{2(b - a)(b - \beta)(\beta - a)}{a^2(a - \beta)^2} - \tfrac{(a - \beta) + a y\eta}{x^2(b - \beta) + a\beta\beta - a\beta}\cdot\tfrac{dy}{dx} = 0,
\]
that is,
\[
\tfrac{dy}{dx} = \tfrac{2(a^2 y + a\eta x)(\beta - a)}{a(a + \beta) - a\eta}\,\tfrac{dy}{dx}.
\]
which is satisfied by the foregoing values of $\tfrac{\delta x}{\delta x}, \tfrac{\delta y}{\delta y}$; the two tangents therefore coincide.

We have
\[
4(by)^2 + 4(by)^2 (4by^2 - 4)\sqrt{a y^2 + b y^2}\,\delta y = \tfrac{4(a + \beta)^2}{a(a + \beta)^2}\,\eta a y - \cdots\xi,
\]
which is virtue of
\[
(ay)(\sqrt{a + \beta})(\beta y^2(\sqrt{a} + a)) + \beta(y^2(\sqrt{a} + a)) + (a(2 + a))y^2 = 4(y\sqrt{a^2 - 4} + 3y^4 y^2(a + \beta)),
\]
is
\[
4[(by)^2(by)^2 + (by^2 + 3y^4 y^2(a + \beta))] = \tfrac{2(a + \beta)^2}{a(a + \beta)^2}\,\eta a y \cdots\xi.
\]
(observe $2(a^2 + \beta) - a(a - \beta) = -(a + (a + \beta)) = -8a^2$ is negative)
\[
= -\tfrac{2a^2}{a(a + \beta)^2}\,\eta a y \cdots\xi.
\]

\clearpage

% --- p.347 ---
\bigskip
\textit{p.~347 [520]}
\smallskip

if $\tfrac{dy}{dx}$ be the value at the entropy. Writing $\tfrac{a^2}{a^2}$ for the element of the are we have
\[
\frac{ds^2}{a^2} = \tfrac{dx^2(a - \beta)^2}{(a - \beta)^2}\!\left[1 - \tfrac{4}{a}\eta\right] - \tfrac{8a^2}{a},
\]
\[
\frac{ds^2}{a^2} = \tfrac{dx^2(-a\beta\gamma)^2}{(a - \beta)^2} - \tfrac{8a^2}{a},
\]
which exhibit the form at the entropy.

\subsubsection*{The Nodal Curve; expressions for the coordinates in terms of a single parameter $a$. Art.\ Nos.\ 53 to 60.}

\textbf{53.} After the foregoing investigation of the nodal curve, I was led to perceive that it is possible to express $\xi, \eta, \xi_0, \eta_0$ in terms of a single variable $a$; and thus to obtain expressions for the coordinates of a point of the nodal curve in terms of the single variable $a$. The analytical form of the nodal curve being real, with its imaginary values of $\xi, \eta$, the modes $y$ may we may assume
\[
\xi = -b^2 - p(b^2 - a^2 - b^2),\quad \eta = -b^2 + p(b^2 - a^2 - b^2),
\]
$(i = \sqrt{-1}$ as usual): viz.\ $p$ denote one or other of the quantities
\[
y_+ = a + (a^2 - b^2),\quad q = a - (a^2 - b^2);
\]
the expressions for $-\beta\gamma x^2$, $-\gamma a y^2$ become
\[
= [(b - p)(a + 1)(b - q)^3](b - p)^3,
\]
and we have therefore the condition that this shall be real (for the two values $b = p$, $\dot{b} = -b$); being real, it will in certain cases be positive, in other cases negative, and the values for the remaining condition is given by the equation
\[
\Delta'(b^2 - \xi^2 - \xi_0^2 - 4)(b^2 + b^2 + \xi^2 + 2b^2)^2 + p^2(b - \xi^2 - 2)(b^2 + \xi)^2 - 0 = 0.
\]

\textbf{54.} The condition of reality is easily found to be
\[
\Delta'(b^2 - \beta\beta - \beta^2)(\beta + \beta - 1) + (a + \beta)^2 + (\gamma + p)^2 + 4(a + p)(\beta - p) + (a + p)^2(b - \beta) - p^2 q^2 = 0;
\]
viz.\ this equation in $\Delta$ must have the form $-a, b$, suppose to a value
\[
\frac{(y - a)\beta}{a^2} = a + \tfrac{a(a^2 - b^2)}{2bg(\gamma - a)}\sqrt{Y}\cdot\eta(a - \beta - \gamma)(a + \beta + \gamma),
\]
where, regarding $X$, $Y$ are given functions of $a$, the coordinates of a plane in a cubic curve having the point $X = 0, Y = 0$ for a node, writing $Y = b + 3a(X + 1)$, $X$ will be such of these is rational function of $a$. The curve has the property that being any line through the entropy, the tangents at the points of the curve, the corresponding entropies on it are
\[
(a^2 + Y)(a^2 + a)(a + 1) + (b + 1)(a^2 + b)\beta(a^2 + a) = -a^2(b - p)^3 + a(p + p)^3,
\]
for a point on the nodal curve.

\textbf{55.} To complete the investigation, writing for a moment $T = 3 = 3a^2(X + 1)$, we obtain
\[
\frac{(b - p)^2}{a^2} = \tfrac{a}{2a^2(a + a - 1)} + \tfrac{a(a^2)(b - \xi)^3}{a + b - 1},
\]
or putting for a moment
\[
\frac{(\gamma - a)\beta}{a^2} = K,
\]

\clearpage

% --- p.348 ---
\bigskip
\textit{p.~348 [520]}
\smallskip

% TODO: book page 348 (algebraic intricate page)
and writing $P + Q = P, X. P = P'$, the first of these is
\[
\tfrac{(y - p)^2}{-\beta\gamma}\beta + \tfrac{X(X + 2)}{(X + 1)^2}\,K - \tfrac{Y(X + 1)\beta}{2(X + 1)} + \tfrac{1 - X}{Y + X},
\]
which regarding $X$, $Y$ as the coordinates of a point in a plane in a cubic curve having the point $X = 0, X = 0$ for a node: hence writing $Y = b + 3a(X + 1)$, $Y$ will be such of these is rational function of $a$. The curve has the property that being any line through the entropy, the tangents at the points of the curve, the corresponding entropies on it are
\[
(a^2 + Y)(a^2 + a)(a + 1) + (b + 1)(a^2 + b)\beta(a^2 + a) = -a^2(b - p)^3 + a(p + p)^3,
\]
for a point on the nodal curve.

\textbf{55.} To complete the investigation, writing for a moment $T = 3 = 3a^2(X + 1)$, we obtain
\[
\frac{(b - p)^2}{a^2} = \tfrac{a}{2a^2(a + a - 1)} + \tfrac{a(a^2)(b - \xi)^3}{a + b - 1},
\]
or putting for a moment
\[
\frac{(\gamma - a)\beta}{a^2} = K,
\]

\clearpage

% --- p.349 ---
\bigskip
\textit{p.~349 [520]}
\smallskip

we have
\[
X = \tfrac{a(K - 7)(a - 1)}{a + 1 - K},\quad X_0 = \tfrac{K(a - 7)(a + 1)}{a + 1 - K},
\]
\[
Y = K \tfrac{3K(a^2 - 1) + 2(a + 1)}{a + 1 - K},\quad Y_0 = K \tfrac{3(a - 1)(K - 2)(a + 1)}{a + 1 - K},
\]
or substituting for $K$ its value we have
\[
K(a - 1) = \tfrac{(a - 1)^2}{ab}\!\left(a^2 - \tfrac{8aq}{(a - 1)^2}\right) = -\tfrac{(a - 1)^2}{ab}\!\left(a + \tfrac{2q}{a - 1}\right)\!\left(a - \tfrac{2q}{a - 1}\right),
\]
$a a + 1 - K = -\tfrac{1}{ab}(a^2 + 1)\!\left(a + b - 1\right)(a - p), \cdots$
$+ b\!(a - p)$,
if as before $\Omega = \beta - \gamma$, and therefore changing the sign of the radical we have the values of $\tfrac{x}{a}, \cdots$,

\textbf{56.} Write for a moment,
\[
\left[\Delta - \tfrac{1}{2}(\gamma - a)\sqrt{1 + \sqrt{\tfrac{2a(a - 2 - 4\beta + 2)}{a(a - 2)}}}\right] - (a - 8a\beta + \sqrt{\beta})Y + B\sqrt{\delta} = 0,
\]
\[
\left[\Delta - \tfrac{1}{2}(\gamma - a)\sqrt{1 + \sqrt{\tfrac{2a(a - 2 - 4\beta + 2)}{a(a - 2)}}}\right] - (a - 8a\beta + \sqrt{\beta})Y + B\sqrt{\delta} = 0,
\]
then in the product of these two expressions the rational part is $= A^2A + B^2B$; and from the manner in which they were arrived at we have $9 = A_1A + B_1B$, and the rational part is thus
\[
= -\tfrac{16}{a^2}(B^2 - B'B).
\]

\clearpage

% --- p.350 ---
\bigskip
\textit{p.~350 [520]}
\smallskip

We have
\[
A^2A - B^2B = (\Delta - a)^2 X - 8a^2 X_0,
\]
\[
B = \tfrac{1}{2}(\gamma + a)X\!\left[a - \tfrac{2aq}{(a - 1)^2}(X + S)\right] - a^2 b^2(P^2 - 2a^2),
\]
hence the rational part in question is
\[
= \pm \tfrac{(\eta + a)^2}{(2a + Pp + a)}\sqrt{\tfrac{a + R + (1 - \xi)(2 + \xi)}{P\{Q + q(P + Q) - 3\xi_0\}}}(2 + \xi).
\]

\textbf{57.} We have
\[
1 = \tfrac{a}{2}\,\tfrac{1}{a^2(a - 1)}\!\left[a a^2 - 2\beta - \xi\!\left(a^2 - \tfrac{4q}{(a^2 - 1)^2}\right)\right]
\]
\[
- \tfrac{1}{2}\,\tfrac{1}{(\gamma - a)^2}\!\left[a a^2 - 2\beta - \xi - \tfrac{4qa}{(a^2 - 1)^2}\right] + \cdots,
\]
where the term in [\ ] is
\[
= \tfrac{2ag}{a^2(a^2 - 1)}\!\left[3a a^2 - 2\beta - \xi\!\left(a^2 - \tfrac{4q\eta}{(a^2 - 1)^2}\right)\right],
\]
hence $\Delta = \tfrac{(\gamma - a)\beta}{(\gamma - 1)}\,\tfrac{x(\beta - a)}{(y - a)} + \tfrac{X(\beta - y)}{(y - a)} - 2y,$
where now
\[
a = \tfrac{1}{2}(P + Q - 1)(\beta - q).
\]

\textbf{58.} Moreover
\[
(\Delta - a)^2 = -\Delta^2 - \Delta a + a^2,
\]
where the term in $[\ ]$ is
\[
= -\tfrac{2ag}{a^2(a^2 - 1)}\!\left[3a a^2 - 2\beta - \xi\!\left(a^2 - \tfrac{4q\eta}{(a^2 - 1)^2}\right)\right];
\]
where on changing $a$ into $\beta$, the coefficient of $a$ is
\[
= \tfrac{2bx}{x^2}\,(b - 1)^2 = -\tfrac{ab(a - b - 1)(a + b)}{x^2}\!\left[3a(a - 2\beta - 4q\beta + 8\beta(a - q + 1) + 8\beta(a^2 - 1)\right]
\]
\[
- \tfrac{a^2(b - a)}{(a^2 - 1)}\!\left[3a a^2 + 4(b - \beta) + 8(\beta - q) + 8(a^2 - 1)\right] + \tfrac{a(b + a)}{a^2(a - 1)}\Sigma,
\]
where the last expression on a single fraction is
\[
= -\tfrac{8ax}{a(\beta - 1)^2}\,(b\beta - q)(a - q)\cdots
\]
and if we hence determine the values $3(\xi^2 - \xi_0^2) = \eta_0$, we find that the resulting terms divide by $\xi - \xi_0$, and we have
\[
3\xi_0^2(\xi_0 - \xi) - \xi^2[(b + a)\xi + 3\xi_0] = (\xi + \xi_0)\xi_0(b + a)(\xi + \xi_0)(\xi + a)
\]
\[
+ \xi_0 \cdots,
\]
\[
+ \xi_0 \cdots.
\]

Hence observing that by the relation between $\xi, \xi_0$, the first term is
\[
= -3Q\xi_0 + (P + Q)q + 2\xi_0 q,
\]
the equations become
\[
1: \eta : \eta_0 = -[Q + p(P + Q) - 3\xi_0] \cdot Q
\]
\[
: R(2\xi_0 + 2) - \xi^2(P + Q\xi)
\]
\[
: R(3\xi_0^2 - \xi^2) - \xi^2 P(P + Q\xi).
\]

\clearpage

% --- p.351 ---
\bigskip
\textit{p.~351 [520]}
\smallskip

Similarly
\[
\Sigma(\Delta - a)^2 + aX\Sigma a^2(a - 1)\Sigma + a^2(\Sigma + S)
\]
\[
= \tfrac{a^2}{a(a - 1)}\!\left[(b\beta - q)(a - q) + a(a - 1)\Delta - q(b\beta - q)\Sigma\Delta + b\beta(a - q)\Sigma\Sigma - q\beta(a - q)\Sigma\Sigma\right],
\]
where the term in $\{\ \}$ is
\[
= -\tfrac{8\beta}{a^2}\,\eta\Delta\Sigma\eta - \tfrac{a^2(b - 1)}{(\beta - 1)}\beta\beta a\Sigma + 3\beta\Sigma(a - 1)\eta - \alpha + \beta - q(\beta - q) - \beta,
\]
in which the coefficient of $a$ is $a\beta(\beta - q)\Sigma\Sigma\Sigma + 3y\eta(a - q)$, but the last is a product of these linear functions: hence
\[
3(\Sigma - aB)^2 + 6\Sigma = \tfrac{2(a - 1)}{a^2}\,\eta(\Sigma - q + 3\beta(a - q) - \delta(a - q) - a).
\]

\textbf{59.} Substituting these values we have the expression
\[
\tfrac{1}{(\beta - 1)\eta\alpha}\!\left[\tfrac{(y - a)q + a(a - 1)(b\beta - q) + 3a\Sigma(a - q) - 2bq + (b\beta - q)^2 - (a\Sigma)}{a\eta(\eta - a)\Sigma a(a)}\right] + 5\Sigma^2,
\]
which solving therein $\Delta = -y - \beta\beta\eta y - 5\beta\eta - 5\beta\eta + \Sigma\Sigma - 5\beta\eta\eta + 4(\eta - q + 1)\Sigma\eta + \cdots$, and the final results are
\[
-\beta\beta z^2 = \tfrac{(y - a)+ \delta(b\beta - q)\Sigma\Sigma + \tfrac{1}{2}\beta(a\Sigma) + \cdots(\Sigma + q^2(\Sigma + 5)(a - 1)}{4y(y - 1)(\beta - 4)\Sigma\Sigma(a - q)},
\]
which is in be observed that, equating the denominators by $\Sigma\Sigma$, we have a simple root $\Sigma\Sigma$, to indicate this, we may insert in the denominator the factor $(1 - 8\beta)$.

\textbf{60.} We see here the meaning of all the factors, viz.\ \ \ \ \ \ \ \ \ \ \ Planes.
\begin{center}
\begin{tabular}{l|cccc}
 & $a = 0$ & $\beta = 0$ & $c = 0$ & $a = \infty$\\
\hline
Evolute nodes & $a = -\tfrac{a^2}{y - a}$ & $a = -1$ & $\cdots$ & $a = -\infty$\\
Umbilicar centres & $a = -\tfrac{2y\beta}{(\eta - \beta)^2}$ & $\cdots$ & $\cdots$ & $\cdots$\\
Entropies & $a = -\tfrac{a + \beta}{a + \beta}(\eta - \beta\beta)$ & $\cdots$ & $\cdots$ & $\cdots$\\
\end{tabular}
\end{center}

\clearpage

% --- p.352 ---
\bigskip
\textit{p.~352 [520]}
\smallskip

For the real curve $a$ extends from $\tfrac{a}{2b}$ through 0 to $-\tfrac{3a}{2b}$, viz.
\[
a = -\tfrac{2y}{y - a},\quad a = -1,\quad \text{gives entropy in plane $a = 0$,}
\]
\[
a = \tfrac{2a\beta}{y - a},\quad \text{umbilicar centre in plane $y = 0$,}
\]
\[
a = -\tfrac{a + \beta}{a + \beta}(\eta - \beta),\quad \text{evolute-node in plane $a$.}
\]

It is to be noticed that the order of magnitude of the terms in the table is
\[
a, -\tfrac{2y}{y - a},\quad a, \tfrac{2a\beta}{y - a},\quad a, \cdots,
\]
$\xi \cdot \tfrac{2y - a + \beta - \tfrac{a + \beta}{a + \beta}}{(\eta + a + \beta)}$; which belong to the real curve are contiguous; this is in fact be should be: Several of the proving investigations conducted by means of this quantities $\xi, \xi, \xi_0$, might have been conducted more simply by means of the formula involving $a$.

\subsubsection*{The Eight Cuspidal Conics. Art.\ Nos.\ 61 to 71.}

\textbf{61.} The centro-surface is the envelope of the quadric
\[
\frac{a^2 x^2}{(a + \xi)^2} + \frac{b^2 y^2}{(b + \xi)^2} + \frac{c^2 z^2}{(c + \xi)^2} - 1 = 0.
\]
Hence it has a cuspidal curve given by means of this equation and the first and second derived equations
\[
\frac{a^2 x^2}{(a + \xi)^3} + \frac{b^2 y^2}{(b + \xi)^3} + \frac{c^2 z^2}{(c + \xi)^3} = 0,
\]
\[
\frac{a^2 x^2}{(a + \xi)^4} + \frac{b^2 y^2}{(b + \xi)^4} + \frac{c^2 z^2}{(c + \xi)^4} = 0,
\]
which equations determine $\alpha^2$, $b^2 y^2$, $c^2 z^2$ in terms of $\xi$, viz.\ we have
\[
-\beta\gamma a^2 x^2 = a^2(a + \xi)^3,
\]
\[
-\gamma a b^2 y^2 = b^2(b + \xi)^3,
\]
\[
-a\beta c^2 z^2 = c^2(c + \xi)^3;
\]
so that, comparing with the equations $-\beta\gamma x^2 = a^2(a + \xi)^3(a + \eta)$, etc., which gave the centro-surface, we see that the cuspidal curve $\xi = \eta$, the cuspidal curve must be in question arises from the entire case in the ellipsoid, which form as the envelope of the curves of curvature: it is clear that the curve is imaginary.

\clearpage

% --- p.353 ---
\bigskip
\textit{p.~353 [520]}
\smallskip

\textbf{62.} From the foregoing equations we have
\[
\sqrt{a x} + \sqrt{\beta}y + \sqrt{c}z = a + \sqrt{\xi},
\]
\[
a\sqrt{a x} + \beta\sqrt{\beta}y + c\sqrt{c}z = a + \sqrt{\xi}\,\xi;
\]
the second of which is best written in the rationalised form
\[
(1, 1, 1, -1, -1, -1)(a x, b y, c z; \beta\sqrt{\gamma}y, \sqrt{\gamma a}z, \sqrt{a\beta}y)^2 = a y^2,
\]
and combining herewith the equation
\[
\sqrt{a x} + \sqrt{\beta}y + \sqrt{c}z + \sqrt{\xi} = 0,
\]
then for any given signs of $\sqrt{a}, \sqrt{\beta}, \sqrt{c}, \sqrt{\xi}$ the first of these equations represents a quartic surface, the second a plane; or the two equations together represent a conic.

By changing the signs of the radicals (observing that when all the four signs are changed simultaneously the curve is unaltered) we obtain in all 8 conics, but only four quartic surfaces; viz.\ the two conics
\[
\sqrt{a x} + \sqrt{\beta}y + \sqrt{c}z = \sqrt{\xi},
\]
\[
\sqrt{a x} + \sqrt{\beta}y + \sqrt{c}z = -\sqrt{\xi}
\]
lie on the same quartic surface.

\textbf{63.} The conics form two tetrads, and the principal conics (reckoning as one of them the conic at infinity) another tetrad: the complete cuspidal curve consists therefore of three tetrads of conics; with these we may form (one or eight out of each tetrad) two tetrads of skew hexagons (omit\.); the corresponding sides of these tetrads, and with a determination omitted in regard to the ellipsoid, the cuspidal conic of the principal eight conic of the principal eight; the same quadric surface. It is in be remembered that, having in mind the corresponding multilinear conics, and equally that the cuspidal conics three of the evolute at that point, which point is also a node on the surface.

\textbf{64.} The 8 conics have two tetrads, and the principal conics (reckoning as one of them the conic at infinity) another tetrad: the complete cuspidal curve consists therefore of three tetrads of conics; with these we may form (one or eight out of each tetrad) two tetrads of skew hexagons (omit\.); the corresponding sides of these tetrads, and with a determination omitted in regard to the ellipsoid, the cuspidal conic of the principal eight conic of the principal eight; the same quadric surface. It is in be remembered that, having in mind the corresponding multilinear conics, and equally that the cuspidal conics three of the evolute at that point, which point is also a node on the surface.

\clearpage

% --- p.354 ---
\bigskip
\textit{p.~354 [520]}
\smallskip

\textbf{65.} But first consider the two conics
\[
\sqrt{a}x + \sqrt{b}y + \sqrt{c}z = \sqrt{-a\beta\gamma},
\]
\[
(1, 1, 1, -1, -1, -1)(a x, b y, c z; \beta\sqrt{\gamma}y, \sqrt{\gamma a}z, \sqrt{a\beta}y)^2 = a y^2;
\]
for the intersections with the plane $y = 0$ we have
\[
\sqrt{a}x + \sqrt{b}y + \sqrt{c}z = \sqrt{-a\beta\gamma},
\]
\[
(a x x - a^2 x - \gamma y z)^2 = 0;
\]
so that the two conics each meet the plane in question in the same two coincident points, that is, they each touch the plane $y = 0$ at the same point. The two points by the equations
\[
\sqrt{a}x + \sqrt{c}z = \sqrt{-a\beta\gamma},
\]
\[
(a x x - y z)^2 = 0;
\]
viz.\ this is the point, $a a = \tfrac{-a y c}{a - \beta}$, $a a = \tfrac{2y\beta}{a - \beta}$, which is one of the umbilicar centres before mentioned. We have
\[
(a^2 + a^2)\,\cdots,
\]
and the common tangent at this point is
\[
\sqrt{a}x + \sqrt{b}y + \sqrt{c}z = \sqrt{-a\beta\gamma},
\]
which is also the common tangent of the ellipse and evolute in the plane $y = 0$.

\textbf{66.} It has been seen that the nodal curve meets each principal conic at four umbilicar centres, which form ten cusps of the nodal curve. It is to be further shown that the nodal curve meets each of the 8 cuspidal conics in the four points (giving 32 nodal points distinguished as the principal entropy) or 16 entropies, and the points now in question are cusps of the nodal curve.

\textbf{67.} Putting for shortness,
\[
\Theta = \sqrt{1\!\left(\!\sqrt{\eta - q}(\eta - \beta)\,(\eta - 2)\right)},
\]
\[
S = \tfrac{a(a + \beta a)}{\eta(\eta - q)}\!\left(\eta - \tfrac{2y}{\eta - a}\right),
\]
we have then
\[
\Theta(1 + a\Theta)\Delta = -\sqrt{a x}.
\]

\clearpage

% --- p.355 ---
\bigskip
\textit{p.~355 [520]}
\smallskip

or, what is the same thing,
\[
\Theta(1 + 3S) - 1 + a\,\sqrt{\Sigma}(2 + S - 2) = 0.
\]

we have without difficulty
\[
\Theta(1 + S) + 1 = \tfrac{1}{2(\eta - q)}[3a^2 - 4a + 4 + \tfrac{8\eta q}{(\eta - q)^2}],
\]
\[
\Theta(1 + 2 + \Theta) - 1 = \tfrac{1}{2(\eta - q)}\!\left[\tfrac{2(\eta - 1)\beta + 2y(\eta + \beta - q)}{(\eta - q)} - \tfrac{4(\beta q + \beta\eta)}{(\eta - q)^2}\right];
\]

Omitting it, the equation becomes
\[
\sqrt{a^2 - \beta - \tfrac{2y}{(\eta - q)\eta}}\sqrt{a^2 - \tfrac{(2)\eta - 4\beta + \tfrac{1}{2}}{(\eta - q)}}\!\left(\eta - 4\eta + 4\eta + \tfrac{2}{(\eta - q)}\right) = 0,
\]
or putting for shortness
\[
\sqrt{a^2 - \beta - \tfrac{2y}{(\eta - q)\eta}} = M,\text{ and rationalising, this is}
\]
\[
= (\sigma - 2 - M)(\Sigma\eta - 2) + 4y\Sigma + 4(\Sigma - 1)
\]
and, working this out, the terms in $\Sigma$ disappear, and we have for the resulting terms divide by $\Sigma - 2$, and we have
\[
8 : a = \cdots = -P\Sigma\beta + 4(\Sigma + \beta + 3\Sigma - 4)\Sigma + 8M,
\]
a quartic equation in $a$: each of the 4 roots there correspond 8 intersections, viz.\ there will be in all 32 intersections, lying in 4's upon the 8 cuspidal conics.

\textbf{68.} To start from these points are cusps, or stationary points on the nodal curve, starting from the expression of $-\beta\gamma z^3$ etc. in terms of $a$ we have, first for $\delta x$, the equations
\[
\tfrac{\delta x}{x} = \tfrac{a}{a - \beta + 1}\,\delta a + \tfrac{4}{a + 1}\,\delta a + \tfrac{2(\eta - \beta)}{(a + \beta + a)(\sqrt{\Sigma} - q + \beta + a)}\,\delta a,
\]
so, as this may be written,
\[
\tfrac{\delta x}{x} = \tfrac{a(2a + 4(a + 1)\beta - q(a + 1) - \beta)\eta(\Sigma + M + 2a - 2P) + (M + 3\eta)\eta - 4a + M)}{a(a + 1)(\beta - 1)\Sigma a\eta - q(a)}
\]
\[
- \tfrac{a\Sigma\,[5 + 5\Sigma\,(\tfrac{1}{2}(a + 8\eta + 9\Sigma) + \tfrac{1}{2}(3M - 8\Sigma\eta\Sigma) + \tfrac{1}{2}(3M(a + 4\eta + a) + \Sigma - 2\Sigma a + a + \beta))]}{a(a + 1)(\Sigma\eta - q)\Sigma\Sigma a(a - q)}
\]
viz.\ the numerator vanishes when $a$ is a root of the quartic equation.

\clearpage

% --- p.356 ---
\bigskip
\textit{p.~356 [520]}
\smallskip

\textbf{69.} We have moreover
\[
-\tfrac{2\delta y}{y}\,da = \tfrac{1}{\eta\eta(\eta - q)}\!\left[a + \tfrac{2\eta}{(\eta - a)(\eta - q)} + \tfrac{2a(\eta - q + a)}{(\eta - a)} + \tfrac{2(a - q)}{(\eta - a)} - \Omega\,\tfrac{6}{a}\right],
\]
which, putting $\tfrac{\eta - q + a}{\eta - q} = R$, and therefore $\tfrac{a - q}{\eta - q} = R - 1$, $a = \tfrac{\Omega}{\eta - q} - C$, is
\[
= \tfrac{1}{2}\!\left[\tfrac{1}{a + R} + \tfrac{1}{R} - 2R + \tfrac{2R}{\eta - q}\cdot\tfrac{4}{\eta - q}\cdot\tfrac{M + 8\eta\beta}{q} + \tfrac{6}{\eta - q}\right];
\]
and adding the fractions inside of $\{\ \}$, the numerator is
\[
a(P R + 5G\Omega - R) + \cdots
\]
\[
= \tfrac{2}{\eta}\,a(\eta + a + 1)(\eta + 5G + 7\eta + 3G\beta) - 6G + (\eta + 5G + 7\eta + 4G)\Omega - 96b\beta\beta + 8G,
\]
which observing that $R = 2 + a$, is $= 8M + 16\beta\beta$,
and, substituting for $M$ and $R$ their values, this is
\[
= a^2(7M + 5R) + (3M + 8\beta) + (\beta + 5G + 2\eta + a) + 8\beta M,
\]
and substituting, the whole denominator is $1 \cdot 8M$, so that, as the result of the equation
\[
8M = a^4 - 3M(\eta + 5G + 2\eta + a) - 8\beta M,
\]
we may have
\[
-\tfrac{2\delta y}{y}\,da = \tfrac{4(2\eta - 1)(\eta + a + 1)}{(\eta - q)\eta}\!\left[a + \tfrac{2\eta}{(\eta - q)}\right] + \tfrac{a}{\eta + a + 2}.
\]
and the numerator of the expression on a single fraction is
\[
= \tfrac{4(2\eta + a)^2}{(\eta - q)^2 \eta^2}\,a(\eta + \tfrac{2y\Sigma}{\eta(\eta - q)}\,a + a)\,\beta(\eta - 4M)\,\cdots,
\]
which is
\[
= 5(\eta + 5G + \eta + \beta) + 8\Omega(\eta + a M + 8\beta a + a) - 24\Omega(\eta + a M + 5\beta\eta + a M),
\]
which is
\[
= a\,\beta = -3M(\eta + 3G + \eta + 8\beta + a + 5\beta) + \tfrac{4(2\eta + a)}{(\eta - q)^2\,\eta}\,a(3M + 5G + 7\eta + 4G);
\]

\clearpage

% --- p.357 ---
\bigskip
\textit{p.~357 [520]}
\smallskip

the term in [\ ] is
\[
= a^2(7M + 5R) + B + 8\,\tfrac{a + \beta a}{(\eta - q)\eta}\,\sigma
\]
\[
= a + 2(M + 8R) + 8\,\tfrac{a + \beta\eta}{(\eta - q)\eta}\!\left(a - 8a + \beta a + a\right)
\]
\[
+ 2[3M\beta + 8\eta M\,\tfrac{a + \beta\eta}{(\eta - q)\eta}\!\left(P + 2\beta\beta + 4G\Omega\eta\right)
\]
\[
- 8M = a\sqrt{a^2(11M + 5R) + 5R(a^2 - 5R) + 5\Sigma + 5R}\,\beta + 5\Omega M],
\]
which is found to be
\[
= -8a\,\Omega + a^2(11M + 5R) + B\cdot\beta + \beta M;
\]
and the whole numerator is
\[
= 8a\,\Omega^2 + a^2(11M + 5R) + B\cdot\beta + \beta M,
\]
which is
\[
= (16 + 5R)\,a + 5R(a^2 - 5R) + 5M + 4P\beta + 4M.
\]

\textbf{71.} We have then
\[
-\tfrac{2\delta y}{y\,da} = \tfrac{1\eta(R + 6\Sigma\beta)(\eta + a) - \tfrac{5G + 2\eta + a}{(\eta - q)^2}\,a + 5RM - 16M^2(\beta + R)\eta + \tfrac{1\eta}{a + 1\beta}\,\beta + 16M\eta + 5R}{(\eta - q)^2\,\eta\beta(\eta - q)},
\]
and theme also
\[
\tfrac{2\delta y}{y\,da}\!\!\left(\tfrac{a + \beta}{\eta - q}\!\right)\Omega + \tfrac{(2\eta + a)(2\eta + a)(\eta + a + a + 1)}{(\eta - q)\,\eta\beta}\,\sigma + \tfrac{a + \beta}{(\eta - q)}\!\left(a\,\tfrac{a}{a - q}\right)
\]
so that $\delta x$ do also vanish when $a$ is a root of the quartic equation, the points in question are therefore cusps of the nodal curve.

\subsubsection*{Centro-surface as the envelope of quartic $\Sigma a^2 = \beta + 1 = 0$. Art.\ Nos.\ 72 to 76.}

\textbf{72.} The equations $-\beta\gamma x^2 = a^2(a + \xi)^3(a + \eta)$, etc., considering therein $\xi, \eta$ as variable parameters connected by the equation $\xi - \eta$ as a given constant value, give the equation of the sequential centro-curve, and considering $\xi$ as a given constant but $\eta$ as variable they give the constant centro-curve.

\textbf{73.} Suppose first that $\eta$ is a given constant; to eliminate $\xi$ we may write the equations in the form
\[
-(\beta\gamma\sqrt{a} = a^2(a + \xi)^{3/2}(a + \eta)^{1/2},\text{ etc.,}
\]
and then multiplying first by $a^{-3/2}, b^{-3/2}, c^{-3/2}$, and adding, and secondly by $b, c, a$, and adding, (observing that $\Sigma\!\left(\sqrt{a}\right) = 0$, $\Sigma(\sqrt{a}\sqrt{b}) = 0$), we obtain
\[
\Sigma(a x)\sqrt{a^2 + a + b} = -a\sqrt{a + b},
\]
\[
\Sigma(\beta a x)\sqrt{a^2 + a + b} = -a(a + b\,\sqrt{a + \eta}\,\beta),
\]
which equations, considering therein $a$ as a given constant, are the equations of a sequential centro-curve.

\clearpage

% --- p.358 ---
\bigskip
\textit{p.~358 [520]}
\smallskip

If from the two equations we eliminate $a$ we should have the equation of the centro-surface; the second equation is the derivative of the first in regard to $\eta$; and it thus appears that the equation of the surface can be obtained by considering $\eta$ in the first equation as a variable parameter. The equation of the centro-surface is thus
\[
\sum (\overline{a x})^2 (a + \eta)^4 = (- a \beta \gamma)^2.
\]

but the form is less convenient to that of \\
The above equation is, multiplying by $b, c, a$, and adding, by $b, c, a$, and adding, by $b^2, c^2, a^2$, and adding,
\[
3a^2 \cdot ab^2 + 3a b^2 c^2 + a b^2 \overline{a x^2} = - a \beta \gamma \cdot \sqrt{(a + \eta) \beta} \cdot ab,
\]
and, eliminating $\eta$, we have the equation in the form before mentioned (last paragraph).

\textbf{74.} Taking now $\xi$ as a given constant, and writing the equation in the form
\[
\sqrt{a} \cdot \overline{a x^2} = a^2 (a + \xi)^{3/2} \cdot (a + \eta)^{1/2}, \text{ etc.,}
\]
then multiplying by $\beta \gamma, \gamma a, a \beta$, and adding, and secondly by $b c, c a, a b$, and adding, or writing this for the sake of distinction
\[
\sqrt{a} \cdot \overline{a x^2} = a^2 \beta^{3/2} (a + \xi)^{1/2}, \text{ etc.,}
\]
we have
\[
\sum (\overline{a x})^2 (a + \xi)^{1/2} = a + \xi,
\]
\[
\sum (\overline{a x})^2 (a + \xi)^{1/2} = a + \xi + b + \xi,
\]
or writing these equations in full form
\[
\sqrt{a^2 (a + \xi)^2 + b^2 (b + \xi)^2 + c^2 (c + \xi)^2 + \delta} - \delta + \xi = 0,
\]
\[
\sqrt{a^2 (a + \xi)^2 + b^2 (b + \xi)^2 + c^2 (c + \xi)^2 + \delta} - 4 + \xi - 1 = 0,
\]
thus the cofficients on the second side are $(a, x, y, z, ax, by, cz, a^2)$, the first quadrilateral. The second equation as the derivative of the first in regard to $\xi$, $a$, $z$, etc., as the first quadrilateral. The second equation as the derivative of the second in regard to $\xi$, $a$, $z$, etc., as a function of $a$ in this way evidently a twofold factor $\left(1 + \frac{1}{\xi}\right)^2$, viz.\ this is, the centro-surface is the envelope of a sequence of quartic surfaces, each of which lies entirely the centro-surface. The numerical results were not made out, but my own. (see ``On a certain Sextic Form,'' \textit{Camb.\ Phil.\ Trans.\ t.\ xx} (1871), pp.\ 397--523, [486].)

\clearpage

% --- p.359 ---
\bigskip
\textit{p.~359 [520]}
\smallskip

\subsubsection*{Another generation of the Centro-surface. Art.\ Nos.\ 77 to 83.}

\textbf{77.} It may precede, the equation of the centro-surface is obtained as the condition in order that the equation
\[
\sum(a x)\sqrt{a + \eta - 1} = 0
\]
may have equal roots. But taking $a$ as arbitrary constant, this is the derived equation of
\[
\sum a (a x)^2 (a + \eta) = a + \eta + a x,
\]
and as such it will have two equal roots if the last-mentioned equation has the property of expressing that the last-mentioned equation, or what is the same thing, the quartic equation
\[
(a + \eta) + a (a x)^2 + a (b + \eta) + a (b + \eta)^2 + (a (a + \eta) + a (a + \eta))^2 = 0,
\]
has three equal roots. The condition for this are that the discriminant and the subdiscriminant shall each of them vanish, and we have thus the equations to be considered as respectively a quadric and a cubic function of $a$; viz.\ the equations
\[
\langle a, h, h\,\xi \rangle (a, b, c, d, e)^4 = 0
\]
\[
\langle a, h, h\,\xi \rangle (a, b, c, d, e)^4 = 0
\]
where the degrees in $(x, y, z)$ of $a$, $b$, $c$, $d$, $e$ are $0$, $2$, $4$, $6$ and those of $a'$, $b'$, $c'$, $d'$ are $2, 4, 6$, $8$ respectively; the equation of the centro-surface is then
\[
\begin{vmatrix} a, & b, & c, & d, & e \\
a', & b', & c', & d', & e' \end{vmatrix} = 0,
\]
which is of the right order 12; but it would be difficult to obtain thereby the developed equation.

\textbf{78.} For the nodal curve the cubic equation must be satisfied by each root of the quadric equation, or that is the same thing, the quadric function must contain as a factor
\[
\langle a, b, c \rangle \langle a, b', c', d' \rangle = 0,
\]
whose the degrees may be taken as below, viz.
\[
\langle 0, 2, 4 \rangle \langle 0, 2, 4 \rangle = 0.
\]

\clearpage

% --- p.360 ---
\bigskip
\textit{p.~360 [520]}
\smallskip

and the order of the nodal curve is then $= 24$: two of the equations in fact are
\[
\begin{vmatrix} a, & b, & e \\
a', & b', & e' \end{vmatrix} = 0,\quad \begin{vmatrix} a, & b, & e \\
a', & b', & e' \end{vmatrix} = 0,
\]
which are surfaces of the orders 4, 6, etc.; or the nodal curve is a complete intersection $\Sigma$ x 6. By the results above obtained as to the nodal curve, it appears that the two surfaces must have an ordinary contact at each of the 16 umbilicar centres, and a stationary or singular contact at each of the 32 entropies.

\textbf{79.} The derivation of the centro-surface from the surface
\[
\frac{a^2}{a^2 + \xi}\,X^2 + \frac{b^2}{b^2 + \xi}\,Y^2 + \frac{c^2}{c^2 + \xi}\,Z^2 - U^2 = 0
\]
requires to be further explained. The surface, say $V = 0$, is a quadric surface depending on the two parameters $\xi$, $\eta$; the same outside in direction with these of the ellipsoid, and their relative magnitudes obey the equation
\[
\frac{1}{a^2 + \xi}\,X^2 + \frac{1}{b^2 + \xi}\,Y^2 + \frac{1}{c^2 + \xi}\,Z^2 = 1
\]
divided by $a$, $b$, $c$ respectively; viz.\ for the absolute magnitudes observe that the equation may be written
\[
\frac{a^2 - \xi^2}{a^2 + \xi}\,X^2 + \frac{b^2 - \xi^2}{b^2 + \xi}\,Y^2 + \frac{c^2 - \xi^2}{c^2 + \xi}\,Z^2 - 1 = 0,
\]
viz.\ the surface $V = 0$ is a surface through the spheroconic which is the intersection of the conduit surface by the arbitrary sphere $a^2 + b^2 + c^2 + U^2 = 0$: but, while the surface is hereby and by the preceding condition as to the axes completely determined, the geometrical significance is anything but clear.

\textbf{80.} Considering then the quartic surface $V = 0$, depending on the parameters $\xi$, $\eta$, suppose that $a$ remains constant while $\xi$ alone varies; we have then consecutive surfaces $V = 0, V_1 = 0$, which on of these $V_1$ say intersect in a point of the centro-surface; the point in question will depend on the two parameters $(\xi, \eta)$, and if these vary simultaneously we have thus a whole series of points on the centro-surface, but it may simultaneously be that there belong to the locus of points or the centro-surface. To make this evident, we must inquire that, that proceeding by both $a$ being also along the value of $\eta$ values of $a$ which are not for which lines and the variable centres are imaginary by means of points on the centro-surface, it is the quadric superceded as $\xi, \eta$, to satisfy the equation of the centro-surface; thus the form of the line of corresponding to the difference of the nodal curve, but it is shown of curvature for the parameter $\eta$.

\clearpage

% --- p.361 ---
\bigskip
\textit{p.~361 [520]}
\smallskip

the variable parameter be $\xi$, then this is a curve on the 12-fold surface $\Omega = 0$ obtained by the elimination of $\xi$ from the equations $V = 0$, $\delta_\xi V = 0$; viz.\ we have $\delta_\xi V$ as the curve in question on the surface. The equation of the curve is the elimination of $a$ as from the two equations $S = 0$, $T = 0$ gives as above the equation of the centro-surface.

\textbf{81.} The surface $\Omega = \delta - T = 0$ obtained as shown by the elimination of $\xi$ from the equations $V = 0, \delta_\xi V = 0$, (or, what is the same thing, by equating to zero the discriminant of $V$ in regard to $\xi$) may be termed the sociate-surface: we have there this quartic and sociate surfaces $S = 0, T = 0$ intersecting in the before-mentioned curve, which we call the sociate-edge; and the locus of these sociate-edges is the centro-surface.

\textbf{82.} We may if we please, changing the parameter in case of the functions, consider the two series of surfaces $S = 0$, $T = 0$ depending on the parameters $a, a'$ respectively; a surface of the first series will correspond to a curve of the second series for the parameter $a$ equal, say we have two surfaces of equal, $a'$, and have these to a sociate-edge: taking $a, a'$ as the parameters of the two surfaces $a, a$ in nearly equal, we have the locus of these sociate-edges is the centro-surface.

\textbf{83.} We may by what precedes, changing the parameter in case of the functions, consider the two series of surfaces $S = 0$, $T = 0$ depending on the parameters $a, a'$ respectively; a surface of the first series will correspond to a curve of the second series for the parameter $a$ equal, say we have two surfaces of equal, $a'$, and have these to a sociate-edge; taking $a, a'$ as the parameters of the two surfaces $a, a$ in nearly equal, we have the locus of these sociate-edges is the centro-surface.

\textit{a posteriori} verification that the surfaces $V = 0, V' = 0, V'' = 0$ intersect in a point of the centro-surface, can be made without interest; the parameters $\xi_1, \eta$ of the point of intersection being constant to be of the order $\eta$; whence, starting from the variation of $\eta$, we indefinitely small in regard to those arising from the variation of $\xi$, we ought therefore to have identically
\[
\frac{-a(a^2 + \xi_0)\xi'(a^2 + \eta_0)}{a^2 + \xi} = \delta\xi(a^2 + \xi_0) - \eta_0 - \eta - \beta - \delta'
\]
\[
= \frac{a\eta(\xi - \xi_0)\sqrt{(a^2 + \eta_0)}}{(a^2 + \xi)\sqrt{(a^2 + \eta_0)\xi'}}
\]
and by decomposing the right-hand side into its component fractions this is at once seen to be true.

\subsubsection*{Third generation of the Centro-surface. Art.\ Nos.\ 84 and 85.}

\textbf{84.} Instead of the foregoing equation $V = 0$, consider the equation
\[
W = \left(\frac{a^2 x^2}{a^2 + \xi} + \frac{b^2 y^2}{b^2 + \xi} + \frac{c^2 z^2}{c^2 + \xi} + \frac{w^2}{\xi}\right) - \left(\frac{a^2 x^2}{a^2 + \eta} + \frac{b^2 y^2}{b^2 + \eta} + \frac{c^2 z^2}{c^2 + \eta} + \frac{w^2}{\eta}\right) = 0.
\]

\clearpage

% --- p.362 ---
\bigskip
\textit{p.~362 [520]}
\smallskip

The equations $\delta_\xi W = 0, \delta_\eta W = 0$ contain only $\xi$ and $\eta$ respectively; the surface is therefore the envelope of $a^2 - 4xy = 0$; the elimination of $\xi$ from the equations $\delta_\xi W = 0, \delta_\eta W = 0$ would therefore lead to the equation of the centro-surface as before. But the equation $W = 0$ contains a factor of three of the equations $V = 0, \delta_\xi V = 0, \delta_\eta V = 0$, viz.\ to fix the ideas, $W = 0$ is connected with the equation $V = 0$ as above given by the elimination of $\xi$ from the equation $W = 0$ depending only on $\xi$, we have $\xi$ in regard to the value of $\xi$ and in the regard to the value of $\eta$ in the same equation; and we have the same equation $W = 0$ depending only on $\xi$; or showing that the determination of $\xi$, $\eta$ by an equation of $\xi$ form is connected.

\[
(a^2 + \xi)(a^2 + \eta) + \tfrac{\xi y_0^2}{a^2 + \eta} + \tfrac{\xi y_0^2}{(a^2 + \xi)^2(a^2 + \eta)^2} = 1 - 0,
\]
so that the surface $\Omega = 0$ is obtained by eliminating $\xi$ from this equation as the several points of the curve of curvature belonging to the parameter $\eta$; viz.\ regarding $\xi$ as the variable parameter on the ellipsoid, then there is on the centro-surface the curve which is the value of $\xi$ as belongs to the ellipsoid, then this is the curve which is the value of $\xi$ as belongs to the ellipsoid; or, in a like manner, to the locus of the sociate edges is the centro-surface.

\textbf{85.} As $\xi$ is constant the curve of of the centro-surface for points moved at distance $a^2$ in the curve, viz.\ in the same way, but as on the cone $a$ may be reduced to a twofold form $\left(1 + \tfrac{1}{\xi}\right)^2$, if $w^2 = 0$, and it has evidently a twofold factor on each side of the various contains the planes (when $a$ has any other value), but it must necessarily form quartic surface envelope of the sphere (whose other generator $\xi$) which has the points (one or eight equal to the centre).

\subsubsection*{Reciprocal Surface. Art.\ No.\ 86.}

\textbf{86.} The centro-surface is the envelope of
\[
\frac{a^2 X^2}{a^2 + \xi} + \frac{b^2 Y^2}{b^2 + \xi} + \frac{c^2 Z^2}{c^2 + \xi} - W^2 = 0;
\]
hence the reciprocal surface is the envelope of $a^2 + y^2 + z^2 + w^2 = 0$ in regard to
\[
\frac{(a^2 + \xi)\,x^2}{a^2}+\frac{(b^2 + \xi)\,y^2}{b^2}+\frac{(c^2 + \xi)\,z^2}{c^2} - w^2 = 0,
\]
that is,
\[
a^2 x^2 + b^2 y^2 + c^2 z^2 + a^2 y^2 - 0,
\]
viz.\ the envelope is
\[
(a^2 x^2 + b^2 y^2 + c^2 z^2)(x^2 + y^2 + z^2) = w^2 (a^2 + b^2 + c^2) - x^2 + y^2 + z^2 - X^2 = 0,
\]
or, expanding and multiplying by $a^2 b^2 c^2$, this is
\[
a^2 b^2 c^2 (a^2 + y^2 + z^2)(a^2 X^2 + b^2 Y^2 + c^2 Z^2) = w^4 (a^2 X^2 + b^2 Y^2 + c^2 Z^2) - a^2 b^2 c^2 (a^2 X^2 + b^2 Y^2 + c^2 Z^2)\,w^2,
\]
or, what is the same thing,
\[
a^2 b^2 c^2 \,V_X X^2 + b^2 c^2 a^2 \,V_Y Y^2 + c^2 a^2 b^2 \,V_Z Z^2 \,X^2 + c^2 a^2 b^2 \,V_W W^2 + a^2 b^2 c^2 + \cdots = 0,
\]
which may be written,
\[
a^2 X(a^2 b^2 + b^2 c^2 + c^2 a^2) \cdot X^2 + (a^2 b^2 c^2 - X^2)\,a^2 W^2 = 0,
\]
that is,
\[
(a, b, c, 1)\,(X^2, Y^2, Z^2, W^2)^2 = 0,
\]
and, consequently,
\[
a^2 + b^2 + c^2 - W^2 - a^2 b^2 c^2 + (a^2 b^2 + b^2 c^2 + c^2 a^2) - a^2 b^2 c^2 = 0,
\]
It would doubtless be interesting to discuss this surface as it here presents itself,
and with reference to the geometrical significance as the locus of the pole, in regard to the sphere, of the plane through the intersecting consecutive normals of the ellipsoid; but I abstain from any consideration of the question.

\subsubsection*{Delineation of the centro-surface for given numerical values of the constants. Art.\ Nos.\ 87 and 88.}

\textbf{87.} I constructed on a large scale a drawing of the centro-surface for the values
\[
a = 50,\quad b = 15,\quad c = 13.
\]

(These were chosen so that $a, b, c$ should have approximately the integer values 7, 3, 4, and that $a^2 + b^2$ should be will greater than $2b^2$; they give a good form of surface,

\clearpage

% --- p.364 ---
\bigskip
\textit{p.~364 [520]}
\smallskip

though perhaps a better selection might have been made; there is a slight objection to the existence of the relation $a^2 = 25$; as in the projection it brings a cusp of the evolute on the ellipse). We have therefore
\[
a = 50,\quad \beta = -33,\quad \gamma = 25,
\]
the ellipses in the principal planes of the centro-surface are
\[
\frac{x^2}{(2)} + \frac{y^2}{(b^2/a)^2} = 1, \quad \frac{x^2}{(2)} + \frac{z^2}{(b^2/a)^2} = 1,
\]
\[
\frac{x^2}{(216a^2)} + \frac{y^2}{(216a^2)} = 1,
\]
\[
\frac{a x^2}{(216a)^2} + \frac{z^2}{(2a)^2} = 1,
\]
and these determine on each axis the two points which are the cusps of the evolutes. We have moreover for evolutes some $= 2998$, $y = 0$, $z = 1990$; and for the entropy $x = 1127$, $y = 1942$, $a = 0$.

\textbf{88.} For the delineation of the nodal curves (counted partially) we have then to find the values of $\xi, \xi_0$; these are given in terms of $a, a, $ \textit{ante}, No.\ 33 [p.\ 334]; where $a$ is a given function of $a$, and a contains in terms of the constants $a, \beta, \gamma$ and of $a, \beta$. It was thought sufficient to divide the interval into 8 equal parts, that is, the values of $a$ were taken to be $-25, -27.5, \cdots -50.6$. The values of $\xi, \xi_0$, being found, those of $a$ etc. were obtained from these by means of the original equations $(a + \xi)$ of $a$ etc., as ratio. The whole series of the results is given in the annexed Table, and from these the drawing was constructed.

I find also in the neighbourhood of the umbilicar centre $(a = 25 + a)$,
\[
\delta x = -02469q^2,
\]
\[
\delta y = -02469 q,
\]
\[
\delta z = -02191 q^2,
\]
and in the neighbourhood of the entropy $(\xi = 30 \cdot 333 + \xi)$, we have
\[
\delta x = 1127 a,
\]
\[
\delta y = 1764 a,
\]
\[
\delta z = 4082 a^3.
\]

\clearpage

% --- p.365 ---
\bigskip
\textit{p.~365 [520]}
\smallskip

% TODO: book page 365 (large numerical table of values for delineation of nodal curves; not transcribed)

\begin{center}
\textit{[Numerical table of $a$, $\xi$, $\xi_0$, $\eta$, $\eta_0$, $a^2x^2$, $b^2y^2$, $c^2z^2$, $X^2$, $Y^2$, $Z^2$, etc.\ for $a = -25, -27.5, \ldots, -50.6$ --- 12 columns of decimal data omitted]}
\end{center}

The calculations were performed before I had obtained the formula in $a$, which would have given the results more easily.

\clearpage

% ============================================================
% PAPER 521 begins on book p.~366
% ============================================================
\section{[521]}
\subsection{On Dr Wiener's Model of a Cubic Surface with 27 Real Lines; and on the Construction of a Double-Sixer.}

\begin{center}
\textit{[From the Transactions of the Cambridge Philosophical Society, vol.\ \textsc{xii}.\ Part \textsc{i}.\ (1873), pp.\ 366--383. Read May 15, 1871.]}
\end{center}

\subsubsection*{I.}

\textsc{I call} to mind that a cubic surface has upon it in general 27 lines which may be all of them real. We may out of the 27 lines (and this in 36 different ways) select 12 lines forming a ``double-sixer,'' viz.\ denoting six lines by $a_1, a_2, a_3, a_4, a_5, a_6$, and the other six by $b_1, b_2, b_3, b_4, b_5, b_6$, then are these 12 lines such, that any two with the same suffix do not meet (for example $a_1, b_1$ do not meet), but any two with different suffixes (for example $a_1, b_2$ or $a_2, b_3$) meet; that the two lines forming a pair $(a_1, b_1)$, etc., the intersection of the planes which meet contain pairs of these in the figure $(a_2, b_4)$, etc., is wholly outside the figure as in the figures itself, the disposition of the points is as shown in the figure.

The model is formed of plaster, and is contained within a cube, the edge of which is = 16 inches, the equations of these three lines are
\begin{align*}
a_1, & x = 0, \quad y = 0,\\
a_2, & x = 0, \quad z = 1,\\
a_3, & y = 1, \quad x = 1,
\end{align*}
then as two from a meet each other, any two may be both, but it must be a constant each line $A$, except that the two lines of a pair (e.g.\ $a_2, b_1, a_3, b_2$) lie at one and most each other. And such a system of twelve lines leads us at once to construct the sixer (in the remaining fifteen lines (six of them on the surface, the others as the intersections of these which meet, and there are which contains the pairs of lines $(a_1, b_1)$, etc., the intersection of these lines being on the surface, by reason of the consideration. The system of the twelve lines is in considered within a cube, the edge of which is $= 16$ inches, the surface of these lines through the centre of the cube, the axes parallel to the edges of the cube, and the axes of length $= 16$ inches.

\end{document}
